\subsection{主覆叠}

定义 2. —— 设 B 为拓扑空间，G 为群。给 G 赋予离散拓扑。以 B 为基、以 G 为群的主纤维丛称为 B 的以 G 为群的主覆叠。

这一术语是合理的，因为这样的 B-空间是 B 的一个覆叠。

非空拓扑空间上的主 G-丛具有次数，等于 Card G。

命题 4. —— 设 B 为拓扑空间，(E, p) 为 B-空间，G 为从右边作用于 E 的离散拓扑群。下列断言等价：

\begin{enumerate}
\def\labelenumi{(\roman{enumi})}
\item
  B-空间 E 是以 G 为群的主覆叠；
\item
  对每个 \(g \in G\) 和每个 \(x \in E\) ，有 \(p(x \cdot g) = p(x)\) ，映射 p 诱导从 E/G 到 B 上的同胚，且映射 \(\theta: (x, g) \mapsto (x, x \cdot g)\) 是从 E × G 到 \(E \times_{B} E\) 上的同胚；
\item
  群 G 连续且自由地作用于 E，映射 p 是平展的，且其纤维是 G 的轨道。
\end{enumerate}

(i)⇒(ii)：这由 I, 第 91 页以及 I, 第 94 页命题 3 得出。(ii)⇒(iii)：在假设 (ii) 之下，G 的作用法则是连续的，且映射 p 是满射且开映射（TG, III, 第 10 页, 引理 2）。设 \(x \in E\) ；映射 \(\theta\) 诱导从 \(\{x\} \times G\) 到 \(\{x\} \times p^{-1}(p(x))\) 上的双射，故群 G 自由作用，且其轨道是 p 的纤维。若以 e 表示 G 的单位元，则 \(E \times_{B} E\) 的对角线是 \(E \times \{e\}\) 在同胚 \(\theta\) 下的像。由于群 G 是离散的，集合 \(\{e\}\) 在 G 中是开集，从而对角线 \(\Delta_{E}\) 在 \(E \times_{B} E\) 中是开集。由 I, 第 31 页命题 7，映射 p 是平展的。

(iii)⇒(i)：由于 p 的每个纤维都是 G 在 E 上作用的一条轨道，映射 p 是满射。因为 p 是平展的，对 B 的每个点 b，存在 b 的开邻域 U 以及 p 在 U 上的连续截面 s。集合 \(s(\mathrm{U})\) 在 E 中是开集，且 s 诱导从 U 到 \(s(\mathrm{U})\) 上的同胚（I, 第 30 页, 推论 3）。对每个 \(g \in G\)，集合 \(s(\mathrm{U}) \cdot g\) 在 E 中是开集，且这些集合对一切 \(g \in G\) 的并等于 \(p^{-1}(\mathrm{U})\)。若 g 与 \(g'\) 是 G 的两个不同元素，则集合 \(s(\mathrm{U}) \cdot g\) 与 \(s(\mathrm{U}) \cdot g'\) 不相交，因为 G 自由地作用于 E。映射 \(f: U \times G \to p^{-1}(\mathrm{U})\) 由 \(f(u, g) = s(u) \cdot g\) 所定义，对 \((u, g) \in U \times G\) 是连续的，因而是从 \(U \times G\) 到 \(p^{-1}(\mathrm{U})\) 上的与 G 的右作用相容的同胚。按定义，E 因而是 B 的以 G 为群的主覆叠。

例. —— 1) 设 E 为拓扑空间，G 为从右边连续且自由地作用于 E 的离散拓扑群。为使典范映射 \(p \colon \mathrm{E} \to \mathrm{E} / \mathrm{G}\) 使 E 成为 E/G 的主覆叠，充要条件是该映射是平展的（命题 4 的条件 (iii)）。例如，假设存在拓扑空间 X 以及与 G 在 E 上的作用相容的平展映射 \(q \colon \mathrm{E} \to \mathrm{X}\) ；则 E 是 E/G 的主覆叠。事实上，以 \(q' \colon \mathrm{E} / \mathrm{G} \to \mathrm{X}\) 表示由 \(q\) 诱导的映射；我们有 \(q = q' \circ p\)，因此由 I, 第 29 页命题 6 d)，\(p\) 是平展的。

\begin{enumerate}
\def\labelenumi{\arabic{enumi})}
\setcounter{enumi}{1}
\tightlist
\item
  设 \(q \colon \mathrm{E} \to \mathrm{X}\) 为分离平展态射。群 \(\operatorname{Aut}_{\mathrm{X}}(\mathrm{E})\) 赋予离散拓扑后，从左边连续地作用于 E。若空间 E 连通，则此作用是自由的（I, 第 34 页, 推论 2）。以 G 表示群 \(\operatorname{Aut}_{\mathrm{X}}(\mathrm{E})^{\circ}\) ，并给 E 赋予与 \(\operatorname{Aut}_{\mathrm{X}}(\mathrm{E})\) 的作用相反的右作用。由前例，赋予典范映射 \(p \colon \mathrm{E} \to \mathrm{E} / \mathrm{G}\) 的空间 E 是以 G 为群的主覆叠。
\end{enumerate}

\subsection{离散群的逆紧且自由的作用}

定理 1. —— 设 G 为从右边连续作用于拓扑空间 E 的离散群。假设 E 的每个点都有邻域 U，使得对 G 的每个异于单位元的元素 g，有 U ∩ (U · g) = ∅。则典范映射 p: E → E/G 使 E 成为 E/G 的以 G 为群的主覆叠。

依假设，G 在 E 上的作用是自由的，故只需证明映射 p 是平展的（I, 第 98 页, 例 1）。设 x 为 E 的点，U 为 x 的开邻域，使得对 G 的每个异于单位元的元素 g 有 \(U \cap U \cdot g = \varnothing\) 。映射 p 是开的（TG, III, 第 10 页, 引理 2），且诱导从 U 到 \(p(\mathrm{U})\) 的单射、连续且开的映射，因而是同胚，这就证明了映射 p 是平展的。

例. —— 1) 设 n 为整数 \(\geqslant 0\) ，\(\mathbf{P}_{n}(\mathbf{R})\) 为 n 维射影空间（TG, VI, 第 13 页），\(S_{n}\) 为 \(R^{n+1}\) 中的单位球面（TG, VI, 第 9 页）。从 \(S_{n}\) 到 \(\mathbf{P}_{n}(\mathbf{R})\) 上的典范映射使 \(S_{n}\) 成为以 \(\{1,-1\}\) 为群的主覆叠，该群以位似作用；轨道是由直径上相对的点组成的偶。

\begin{enumerate}
\def\labelenumi{\arabic{enumi})}
\setcounter{enumi}{1}
\item
  每个次数为 2 的覆叠都具有唯一的以 Z/2Z 为群的主覆叠结构。
\item
  设 X 为 R 上有限维的局部分离微分流形，\(\widetilde{X}\) 为 X 的切丛的定向流形（VAR, R, 10.2.4）。空间 \(\widetilde{X}\) 连同它到 X 上的典范投影，是以 \(\{1,-1\}\) 为群的 X 的主覆叠。
\end{enumerate}

推论 1. —— 设 G 为从右边连续作用于拓扑空间 E 的离散拓扑群。下列条件等价：

\begin{enumerate}
\def\labelenumi{(\roman{enumi})}
\item
  群 G 逆紧且自由地作用于 E；
\item
  空间 E/G 是豪斯多夫的，且 E 是以 E/G 为基、以 G 为群的主覆叠。
\end{enumerate}

此外，在这些条件之下，空间 E 是豪斯多夫的。

设条件 (i) 成立。则空间 E 与 E/G 都是分离的（TG, III, 第 29 页, 命题 3），且对每个 \(x \in \mathrm{E}\) ，存在 \(x\) 在 E 中的开邻域 U，使得 \(\mathrm{U} \cap (\mathrm{U} \cdot g) = \varnothing\) 对 G 的每个异于单位元的元素 \(g\) 成立（TG, III, 第 32 页, 命题 8）。由定理 1，条件 (ii) 随之成立。

蕴涵 (ii)⇒(i) 由 I, 第 96 页推论 1 得出。

推论 2. —— 设 G 为拓扑群，H 为 G 的离散子群。让 H 以右平移作用于 G。则从 G 到模 H 的左陪集空间 G/H 中的典范映射赋予 G 一个 G/H 的以 H 为群的主覆叠结构。

设 V 为 G 的单位元 \(e\) 的邻域，使得 \(\mathrm{H} \cap \mathrm{V} = \{e\}\) 。存在 \(e\) 在 G 中的开邻域 U，使得 \(\mathrm{U}^{-1} \cdot \mathrm{U} \subset \mathrm{V}\) （TG, III, 第 3 页），从而 \(\mathrm{U} \cap \mathrm{U} \cdot h = \varnothing\) 对所有 \(h \in \mathrm{H}\) （\(h \neq e\) ）成立。因此，I, 第 100 页的推论 2 由定理 1 得出。

例 4. —— 赋予从 R 到 T = R/Z 的典范映射（TG, V, 第 2 页），R 是 R/Z 的以 Z 为群的主覆叠。

注 1. —— 设 G 为分离拓扑群，K 为 G 的紧子群，H 为 G 的离散子群。群 G 从右边连续且逆紧地作用于模 K 的右陪集空间 K\G（TG, III, 第 30 页, 推论）。由于子群 H 在 G 中是闭的，它也逆紧地作用于 K\G（TG, III, 第 27 页, 例 1）。此外，若 \(H \cap gKg^{-1} = \{e\}\) 对所有 \(g \in G\) 成立，则群 H 自由地作用于 K\G。于是由推论 1，空间 K\G 是 K\G/H 的以 H 为群的主覆叠。

推论 3. —— 设 G 与 G' 为拓扑群，\(\varphi\colon G \to G'\) 为连续、开且满的同态。假设 \(\varphi\) 的核 H 是离散的。则对于 H 以右平移在 G 上的作用，\(\varphi\) 使 G 成为 G' 的以 H 为群的主覆叠。

此外，若群 G 连通，则 H 包含于 G 的中心。

同态 \(\varphi\) 诱导从 G/H 到 G' 上的拓扑群同构（TG, III, 第 16 页, 命题 24），由此并利用推论 2 即得第一个断言。设群 G 连通。对每个 \(h \in \mathrm{H}\) ，由 \(g \mapsto ghg^{-1}\) 定义的从 G 到 H 的连续映射是常值的，取值为 \(h\) 。故群 H 包含于 G 的中心。

注. —— 2) 设 G 与 G' 为拓扑群，\(\varphi\colon\mathrm{G}\to\mathrm{G}'\) 为连续且开的同态。若群 G' 连通，则同态 \(\varphi\) 是满射。事实上，\(\varphi(\mathrm{G})\) 是 \(\mathrm{G}'\) 的开子群，从而是闭子群，于是等于 \(\mathrm{G}'\)。

\begin{enumerate}
\def\labelenumi{\arabic{enumi})}
\setcounter{enumi}{2}
\tightlist
\item
  设 G 为无穷可数的局部紧群，G' 为基础空间为贝尔空间的分离拓扑群。每个从 G 到 G' 上的连续满同态都是开的（TG, IX, 第 56 页, 推论及 TG, III, 第 16 页, 命题 24）。
\end{enumerate}

例. —— 5) 对每个整数 \(n > 0\)，以 \(\mu_{n}\) 表示 C 中 n 次单位根组成的群（A, V, 第 75 页）。映射 \(z \mapsto z^{n}\) 使 \(C^{*}\) 成为 \(C^{*}\) 的以 \(\mu_{n}\) 为群的主覆叠，该群在 \(C^{*}\) 中以乘法作用。

\begin{enumerate}
\def\labelenumi{\arabic{enumi})}
\setcounter{enumi}{5}
\tightlist
\item
  赋予映射 \(z \mapsto e^{2\pi iz}\) （TG, VIII, 第 8 页, 注），空间 C 是 \(C^{*}\) 的以 Z 为群的主覆叠。映射 \(x \mapsto e^{2\pi ix} = \mathbf{e}(x)\) （从 R 到 \(S_{1}\) ）使 R 成为 \(S_{1}\) 的以 Z 为群的主覆叠。
\end{enumerate}

\subsection{伽罗瓦覆叠}

定义 3. —— 设 B 为非空拓扑空间。称 B 的覆叠 E 为伽罗瓦覆叠，如果它是连通的，且对 B 的每个点 b，由 E 的 B-自同构组成的群 \(\mathrm{Aut}_{B}(E)\) 在纤维 \(E_{b}\) 上的作用是传递的。

设 E 为拓扑空间 B 的伽罗瓦覆叠，p 为其投影。映射 p 是满射（A, I, 第 56 页, 定义 6），故空间 B 是连通的。因而覆叠 (E, p) 具有次数，且该次数非零。

命题 5. —— 设 B 为非空拓扑空间，E 为 B 的伽罗瓦覆叠。映射 \((h,x)\mapsto h(x)\) （从 \(\mathrm{Aut}_{\mathrm{B}}(\mathrm{E})\times\mathrm{E}\) 到 E 中）是群 \(\mathrm{Aut}_{\mathrm{B}}(\mathrm{E})^{\circ}\) 在 E 上的一个右作用，它使 E 成为以 \(\mathrm{Aut}_{\mathrm{B}}(\mathrm{E})^{\circ}\) 为群的主覆叠。

给群 \(\mathrm{Aut}_{\mathrm{B}}(\mathrm{E})^{\circ}\) 赋予离散拓扑；此时运算法则 \((h,x)\mapsto h(x)\) 是连续的。此作用是自由的（I, 第 34 页, 命题 11 的推论 2）。由于 E 是 B 的伽罗瓦覆叠，其纤维都是该作用的轨道。由 I, 第 97 页命题 4，E 是以 \(\mathrm{Aut}_{\mathrm{B}}(\mathrm{E})^{\circ}\) 为群的主覆叠。

定理 2. —— 设 B 为连通拓扑空间，E 为 B 的连通非空覆叠。下列性质等价：

\begin{enumerate}
\def\labelenumi{(\roman{enumi})}
\item
  覆叠 E 是伽罗瓦的。
\item
  存在离散拓扑群 G 以及 G 在 E 上的连续右作用，使 E 成为以 G 为群的主覆叠。
\item
  E 的覆叠 \(E \times_{B} E\) （由投影 \(\mathrm{pr}_{1}\) 所定义）是可平凡化的。
\end{enumerate}

当这些条件成立时，由 G 的作用所定义的典范映射 \(G \to \mathrm{Aut}_{B}(E)^{\circ}\) 是群同构。

此外，若空间 B 是局部连通的，则上述性质等价于下述性质：

(i') 存在 B 的点 b，使得群 \(\mathrm{Aut}_{\mathrm{B}}(\mathrm{E})\) 在纤维 \(E_{b}\) 上的作用是传递的。

蕴涵 (i)⇒(ii) 由命题 5 得出，蕴涵 (ii)⇒(iii) 由 I, 第 95 页注 1 得出。我们来证明 (iii)⇒(i)。以 \(p\colon E \to B\) 表示覆叠 E 的投影。由于 B 连通，该覆叠具有次数，又因 E 非空，此次数非零。设 \(b\) 为 B 的一点。我们来证明 \(\mathrm{Aut}_{B}(E)\) 在纤维 \(E_{b}\) 上的作用是传递的。由前所述，该纤维非空。设 \(x\) 与 \(x'\) 为 \(E_{b}\) 中的点。B-态射 \(h\colon E \to E\) （满足 \(h(x) = x'\) ）与连续截面 \(s\) 成一一对应，后者是映射 \(\mathrm{pr}_{1}\colon E \times_{B} E \to E\) 的截面，且满足 \(s(x) = (x, x')\) （I, 第 9 页, 命题 3）。在假设 (iii) 之下，这样的截面存在（I, 第 70 页, 推论 2），且因空间 E 连通而是唯一的（I, 第 34 页, 命题 11 的推论 1）。于是，对每个偶 \((x, x') \in E \times_{B} E\)，存在唯一的 B-态射 \(h\colon E \to E\) 使得 \(h(x) = x'\) 。若 \(h'\colon E \to E\) 是满足 \(h'(x') = x\) 的唯一 B-态射，则有 \(h'(h(x)) = x\) 与 \(h(h'(x')) = x'\) ，从而 \(h' \circ h = \text{Id}_E\) 且 \(h \circ h' = \text{Id}_E\) 。这就证明 \(h\) 是 E 的 B-自同构。因而覆叠 E 是伽罗瓦的。

我们已证明条件 (i)、(ii)、(iii) 等价。设它们成立。以 \(\delta\colon\mathbf{G}\to\mathrm{Aut}_{\mathrm{B}}(\mathrm{E})\) 表示把 \(g\in\mathrm{G}\) 对应到 E 的 B-自同构 \(x\mapsto x\cdot g\) 的映射。映射 \(\delta\) 是从 G 到 \(\mathrm{Aut}_{\mathrm{B}}(\mathrm{E})^{\circ}\) 中的群同态。由于 G 自由地作用于 E，映射 \(\delta\) 是单射。设 \(h\in\mathrm{Aut}_{\mathrm{B}}(\mathrm{E})\)，\(x\in\mathrm{E}\) ，并设 \(g\) 为 G 中满足 \(x\cdot g=h(x)\) 的唯一元素。

B-态射 \(h\) 与 \(\delta(g)\) 在点 \(x\) 处取值相同，从而由于空间 E 连通而处处相同（I, 第 34 页, 命题 11 的推论 1），这就证明了 \(\delta\) 是同构。

现在设空间 B 局部连通，我们证明在假设 (i') 之下覆叠 E 是伽罗瓦的。设 \(E'\) 为 E 按群 \(\mathrm{Aut}_{B}(E)^{\circ}\) 的右作用所得的商空间，并以 \(q\colon E \to E'\) 表示典范映射。它使 E 成为 E' 的满覆叠（I, 第 98 页, 例 2）；映射 \(p\colon E \to B\) 过渡到商，定义出连续映射 \(p'\colon E' \to B\) ，满足 \(p'\circ q = p\) 。由于空间 B 局部连通，B-空间 \((E', p')\) 是一个覆叠（I, 第 81 页, 命题 7）。在假设 (i') 之下，存在 B 的点 \(b\) 使得纤维 \(E'_{b}\) 恰有一个元素；由于空间 B 连通，对 B 的每个点 \(b\) 同样成立（I, 第 74 页, 命题 4），这就证明了 E 是 B 的伽罗瓦覆叠。

命题 6. —— 设 B 为拓扑空间，G 为群。设 (E, p) 为 B 的以 G 为群的主覆叠。假设空间 B 连通且局部连通。设 E₀ 为 E 的一个连通分支，G₀ 为 G 中稳定 E₀ 的子群（A, I, 第 51 页）。B-空间 (E₀, p \textbar{} E₀) 对于 G₀ 在 E₀ 上的右动作是一个主覆叠。此覆叠是伽罗瓦的。

由于空间 B 局部连通，E 亦然，故 \(E_{0}\) 是 E 的开子空间（TG, I, 第 85 页, 命题 11）。又因它还是闭的（TG, I, 第 83 页），空间 \(E_{0}\) 因而是 B 的覆叠（I, 第 80 页, 推论 1）。由于 B 连通，该覆叠具有次数；由于 \(E_{0}\) 非空，\(p|E_{0}\) 的一切纤维都非空。设 x 为 \(E_{0}\) 的一点，并设 \(x' \in E_{0}\) 满足 \(p(x') = p(x)\)；存在元素 \(g \in G\) 使得 \(x \cdot g = x'\)。由于连通分支 \(E_{0}\) 与 \(E_{0} \cdot g\) 有公共点，二者相等，由此得 \(g \in G_{0}\)。于是，群 \(G_{0}\) 传递地作用于 B-空间 \(E_{0}\) 在 \(p(x)\) 处的纤维上。命题 6 得证。

注记. —— 若在命题 6 中不假设空间 B 局部连通，则空间 \(E_{0}\) 可能不是 B 的覆叠（I, 第 147 页, 习题 1）。

\subsection{相伴纤维丛}

设 E 与 F 为集合，G 为作用于 E 的右边、作用于 F 的左边的群。群 G 作用于乘积 \(E \times F\) 的右边，其法则为 \((x, y) \cdot g = (x \cdot g, g^{-1} \cdot y)\) ，其中 \(g \in G\)，\((x, y) \in E \times F\) 。 \(E \times F\) 对此作用的商集记作 \(E \times^{G} F\) 。

当 E 与 F 是拓扑空间时，给集合 \(E \times^{G} F\) 赋予由 \(E \times F\) 的拓扑诱导的商拓扑。从 \(E \times F\) 到 \(E \times^{G} F\) 上的典范映射是连续的。若群 G 连续地作用于 E 与 F，则它是开的。

另外，设 \(F'\) 为 G 从左边作用于其上的集合，\(h: F \to F'\) 为与 G 在 F 与 \(F'\) 上的作用相容的映射（A, I, 第 50 页）。映射 \(Id_{E} \times h: E \times F \to E \times F'\) 与 G 的作用相容，并且通过过渡到商，定义一个记作 \(Id_{E} \times^{G} h\) 的、从 \(E \times^{G} F\) 到 \(E \times^{G} F'\) 中的映射。

若 \(h \colon F \to F'\) 是连续映射（与 G 在 F 与 \(F'\) 上的作用相容），则映射 \(\mathrm{Id}_{\mathrm{E}} \times^{\mathrm{G}} h\) 是连续的（TG, I, 第 21 页, 命题 6）。

例. —— 1) 设 F 为拓扑空间，G 为从左边连续作用于 F 的拓扑群。若给拓扑空间 G 赋予 G 以右平移的作用，则空间 G × \(^{G}\) F 典范地等同于 F，如下所述。连续映射 \(\varphi\) : F → G × F 与 \(\psi\) : G × F → F （分别由 \(\varphi(f) = (e, f)\) （其中 e 表示 G 的单位元）与 \(\psi(g, f) = g \cdot f\) 定义）诱导出连续映射 \(\overline{\varphi}\) : F → G × \(^{G}\) F 与 \(\overline{\psi}\) : G × \(^{G}\) F → F，二者互为逆映射。

\begin{enumerate}
\def\labelenumi{\arabic{enumi})}
\setcounter{enumi}{1}
\item
  例 1 可如下推广。设 B 与 F 为拓扑空间，G 为从左边连续作用于 F 的拓扑群。通过过渡到商，由 (b, f) ↦ ((b, e), f) 给出的从 B × F 到 (B × G) × F 的映射，以及由 ((b, g), f) ↦ (b, gf) 给出的从 (B × G) × F 到 B × F 的映射，定义出互为逆映射的 B-同构 B × F → (B × G) × \(^{G}\) F 与 (B × G) × \(^{G}\) F → B × F。
\item
  类似地，设 E 为拓扑空间，G 为从右边连续作用于 E 的拓扑群。若给拓扑空间 G 赋予 G 以左平移的作用，则空间 \(E \times^{G} G\) 等同于 E。
\end{enumerate}

设 E 与 F 为拓扑空间，G 为从右边连续作用于 E、从左边连续作用于 F 的拓扑群。设 B 为拓扑空间，\(p\colon E \to B\) 为满足 \(p(x \cdot g) = p(x)\) （其中 \(x \in E\)，\(g \in G\)）的连续映射。映射 \(p \circ pr_{1}\colon E \times F \to B\) 通过过渡到商，定义一个连续映射 \(p^{F}\colon E \times^{G} F \to B\) ，而典范映射 \(\pi\colon E \times F \to E \times^{G} F\) 是一个 B-态射。

设 B' 为拓扑空间，\(h\colon B'\to B\) 为连续映射。群 G 连续地作用于 \(B'\times_{B}E\) 的右边，其作用法则为 \(((b',x),g)\mapsto(b',x\cdot g)\)。将典范同构 \(((b',x),y)\mapsto(b',(x,y))\) ——它把 \((\mathrm{B}'\times_{\mathrm{B}}\mathrm{E})\times\mathrm{F}\) 映到 \(\mathrm{B}'\times_{\mathrm{B}}(\mathrm{E}\times\mathrm{F})\) 上——与映射 \(\operatorname{Id}_{\mathrm{B}'}\times\pi\) （从 \(\mathrm{B}'\times_{\mathrm{B}}(\mathrm{E}\times\mathrm{F})\) 到 \(\mathrm{B}'\times_{\mathrm{B}}(\mathrm{E}\times^{\mathrm{G}}\mathrm{F})\) 中）复合，得到连续映射 \(\lambda_{0}\colon(\mathrm{B}'\times_{\mathrm{B}}\mathrm{E})\times\mathrm{F}\to\mathrm{B}'\times_{\mathrm{B}}(\mathrm{E}\times^{\mathrm{G}}\mathrm{F})\)。它通过过渡到商，定义一个连续映射

\[
\lambda \colon (\mathrm{B} ^ {\prime} \times_ {\mathrm{B}} \mathrm{E}) \times^ {\mathrm{G}} \mathrm{F} \rightarrow \mathrm{B} ^ {\prime} \times_ {\mathrm{B}} (\mathrm{E} \times^ {\mathrm{G}} \mathrm{F}).
\]

引理. —— 映射 λ 是同胚。

映射 \(\lambda_0\) 是满射，且 \((\mathrm{B}' \times_{\mathrm{B}} \mathrm{E}) \times \mathrm{F}\) 中两个元素在 \(\lambda_0\) 下有同一像，当且仅当它们属于 G 在 \((\mathrm{B}' \times_{\mathrm{B}} \mathrm{E}) \times \mathrm{F}\) 上作用的同一轨道。由此可知映射 \(\lambda\) 是双射。由于映射 \(\pi\) 是开的，映射 \(\mathrm{Id}_{\mathrm{B}'} \times_{\mathrm{B}} \pi\) 也是开的（I, 第 17 页, 命题 8），从而 \(\lambda_0\) 与 \(\lambda\) 也是开的（TG, I, 第 32 页, 命题 3），这就证明了 \(\lambda\) 是同胚。

命题 7. —— 设 B 为拓扑空间，G 为拓扑群，(E,p) 为 B 上以 G 为群的主纤维丛，F 为 G 从左边连续作用于其上的拓扑空间。

\begin{enumerate}
\def\labelenumi{\alph{enumi})}
\item
  拓扑空间 \(E \times^{G} F\) 连同连续映射 \(p^{F}: E \times^{G} F \to B\) （由映射 \(p \circ pr_{1}: E \times F \to B\) 过渡到商而得），是以 F 为纤维型的局部平凡纤维丛；若 B-丛 E 可平凡化，则它也可平凡化。
\item
  设 \(\pi\colon\mathrm{E}\times\mathrm{F}\to\mathrm{E}\times^{\mathrm{G}}\mathrm{F}\) 为典范满射。映射 \(\mu\colon\mathrm{E}\times\mathrm{F}\to\mathrm{E}\times_{\mathrm{B}}(\mathrm{E}\times^{\mathrm{G}}\mathrm{F})\) 使 \((x,\pi(x,f))\) 对应于 \((x,f)\) ，它是一个同胚，其逆映射是 E 上的局部平凡纤维丛 \((\mathrm{E}\times_{\mathrm{B}}(\mathrm{E}\times^{\mathrm{G}}\mathrm{F}),\mathrm{pr}_{1})\) 的一个平凡化。
\end{enumerate}

先设 E 为平凡主纤维丛 \(B \times G\) 。由例 2，此时空间 \(E \times^{G} F\) 等同于 \(B \times F\) ，映射 \(p^{F}\) 等同于第一投影 \(pr_{1}: B \times F \to B\) ，这就证明在此情形 \((\mathrm{E} \times^{\mathrm{G}}\mathrm{F}, p^{\mathrm{F}})\) 是 B 上的可平凡化纤维丛。在一般情形，B 的每个点都有邻域 U，使映射 \(p_{\mathrm{U}}: p^{-1}(\mathrm{U}) \to \mathrm{U}\) 把 \(p^{-1}(\mathrm{U})\) 变为 U 上的可平凡化主纤维丛。由前所述，\((p^{-1}(\mathrm{U}) \times^{\mathrm{G}}\mathrm{F} \to \mathrm{U}, (p_{\mathrm{U}})^{\mathrm{F}})\) 是 U 上的可平凡化纤维丛。把上面的引理应用于 \(B' = U\) 且 \(h: U \to B\) 为典范单射的情形，对 U-空间 \(((p^{\mathrm{F}})^{-1}(\mathrm{U}), (p^{\mathrm{F}})_{\mathrm{U}})\) （由 \((\mathrm{E} \times^{\mathrm{G}}\mathrm{F}, p^{\mathrm{F}})\) 限制到 U 上得到）结论同样成立。这就证明 \((\mathrm{E} \times^{\mathrm{G}}\mathrm{F}, p^{\mathrm{F}})\) 是 B 上的局部平凡纤维丛，从而完成了断言 a) 的证明。

映射 \(\theta\colon\mathrm{E}\times\mathrm{G}\to\mathrm{E}\times_{\mathrm{B}}\mathrm{E}\) 由 \(\theta(x,g)=(x,x\cdot g)\) 所定义，它是同胚（I, 第 94 页, 命题 3），且与 G 在 E \(\times\) G 上和 E \(\times_{\mathrm{B}}\) E 上由 \(((x,g),g')\mapsto(x,gg')\) 与 \(((x,y),g')\mapsto(x,yg')\) 给出的作用相容。因此 \(\theta\) 通过过渡到商，诱导出同胚 \(\theta'\) （从 (E \(\times\) G) \(\times^{\mathrm{G}}\) F 到 (E \(\times_{\mathrm{B}}\) E) \(\times^{\mathrm{G}}\) F 上）。映射 \(\mu\) 是同胚 E \(\times\) F \(\rightarrow\) (E \(\times\) G) \(\times^{\mathrm{G}}\) F （例 2）、\(\theta'\) 以及典范同胚 (E \(\times_{\mathrm{B}}\) E) \(\times^{\mathrm{G}}\) F \(\rightarrow\) E \(\times_{\mathrm{B}}\) (E \(\times^{\mathrm{G}}\) F) （I, 第 105 页, 引理）的复合，由此即得 \(b\) 。

在命题 7 的假设之下，局部平凡纤维丛 (E × \(^{G}\) F, p \(^{F}\) ) 称为与主纤维丛 (E, p) 相伴的、以 F 为纤维型的局部平凡纤维丛。B-空间 E × \(^{G}\) F 的一切纤维都与空间 F 同胚。特别地，若空间 F 是离散的，则映射 p \(^{F}\) 是一个覆叠。

例. —— 4) 设 B 为拓扑空间，G 为拓扑群，(E, p) 为 B 上的主 G-丛。设 H 为 G 的子群。

以 \(\varphi\colon\mathrm{E}\to\mathrm{E}\times(\mathrm{G}/\mathrm{H})\) 与 \(\psi\colon\mathrm{E}\times(\mathrm{G}/\mathrm{H})\to\mathrm{E}/\mathrm{H}\) 表示由 \(\varphi(x)=(x,\mathrm{H})\) 与 \(\psi(x,g\mathrm{H})=(x\cdot g)\mathrm{H}\) 定义的映射。它们与到 B 的投影以及 G 的作用相容，并通过过渡到商，定义出互为逆映射的 B-空间态射 \(\overline{\varphi}\colon\mathrm{E}/\mathrm{H}\to\mathrm{E}\times^{\mathrm{G}}(\mathrm{G}/\mathrm{H})\) 与 \(\overline{\psi}\colon\mathrm{E}\times^{\mathrm{G}}(\mathrm{G}/\mathrm{H})\to\mathrm{E}/\mathrm{H}\) 。我们称 \(\overline{\varphi}\) 为从 E/H 到 \(\mathrm{E}\times^{\mathrm{G}}(\mathrm{G}/\mathrm{H})\) 上的典范同胚。特别地，拓扑空间 E/H 连同连续映射 \(p_{\mathrm{H}}\colon\mathrm{E}/\mathrm{H}\to\mathrm{B}\) ，是 B 上的以 G/H 为纤维型的局部平凡纤维丛。

此外，若 H 是 G 的正规子群，则 G 的作用通过过渡到商，赋予 B-空间 E/H 一个以 G/H 为群的主纤维丛结构。

特别地，若 E 是 B 的以 G 为群的主覆叠，则 E/H 是 B 的覆叠；若 H 在 G 中正规，则它是以 G/H 为群的主覆叠。

\begin{enumerate}
\def\labelenumi{\arabic{enumi})}
\setcounter{enumi}{4}
\item
  设 B 为拓扑空间，G 为拓扑群，E 为 B 上的主 G-丛。设 F 为关于 G 的拓扑齐性空间（TG, III, 第 12 页）。设 y 为 F 的一点，\(G_{y}\) 为其稳定子。映射 \(\varphi_{y}\colon x\mapsto(x,y)\) （从 E 到 \(E\times F\)）通过过渡到商，定义出同胚 \(\overline{\varphi}_{y}\) （从 \(E/G_{y}\) 到 \(E\times^{G}F\) 上）。当群 G 交换时，子群 \(G_{y}\) 不依赖于点 y，但同胚 \(\overline{\varphi}_{y}\) 一般依赖于它。
\item
  设 B 为拓扑空间，G 为拓扑群，\((E,p)\) 为 B 上的主 G-丛。设 H 为拓扑群，\(f:G\to H\) 为拓扑群态射；给 H 赋予由 \(g\cdot h=f(g)h\) 给出的 G 的左作用。
\end{enumerate}

设 \(q \colon \mathrm{E} \times \mathrm{H} \to \mathrm{E} \times^{\mathrm{G}} \mathrm{H}\) 为典范映射；它是开的，从而是普遍严格的（I, 第 20 页, 推论 11）。映射 \(m \colon (x, h, h') \mapsto q(x, hh')\) （从 \(\mathrm{E} \times \mathrm{H} \times \mathrm{H}\) 到 \(\mathrm{E} \times^{\mathrm{G}} \mathrm{H}\) 中）是连续的。设 \((x, h) \in \mathrm{E} \times \mathrm{H}\) ，\(h' \in \mathrm{H}\) 以及 \(g \in \mathrm{G}\) ；我们有 \(m(xg, f(g)^{-1}h, h') = q(xg, f(g)^{-1}hh') = q(x, hh') = m(x, h, h')\) 。因此，存在唯一的映射

\[
m ^ {\prime} \colon (\mathrm{E} \times^ {\mathrm{G}} \mathrm{H}) \times \mathrm{H} \rightarrow \mathrm{E} \times^ {\mathrm{G}} \mathrm{H}
\]

使得 \(m'(q(x,h),h') = q(x,hh')\) 对一切 \(x \in E\) 与一切 \(h, h' \in H\) 成立。由于 q 是普遍严格的，映射 \(m'\) 是连续的。这是拓扑群 H 在 B-空间 \(E \times^{G} H\) 上的一个右作用。

设 U 为 B 的开子集，使得 \(E_{U}\) 同构于 U 上的主纤维丛 \(U \times G\) 。赋予 H 的作用后，U-空间 \((\mathrm{E} \times^{\mathrm{G}}\mathrm{H})_{\mathrm{U}}\) 等同于主纤维丛 \(U \times H\) 。这就证明 \(E \times^{G}H\) 是 B 上的主 H-丛。

定义 4. —— 设 B 为拓扑空间，G 为拓扑群。称 B 上的局部平凡纤维丛 X 与 B 上以 G 为群的主纤维丛 E 相伴，如果存在一个 G 从左边连续作用于其上的拓扑空间 F，以及一个从 E × \(^{G}\) F 到 X 上的 B-同构。

设 E 为 B 上的主 G-丛，X 为 B 上与 E 相伴的局部平凡纤维丛。若主丛 E 可平凡化，则 X 可平凡化（命题 7）。若 B' 为拓扑空间且 \(h \colon \mathrm{B}' \to \mathrm{B}\) 为连续映射，则由 X 经换基得到的局部平凡纤维丛 \(\mathrm{B}' \times_{\mathrm{B}} \mathrm{X}\) 与 B' 上的主纤维丛 \(\mathrm{B}' \times_{\mathrm{B}} \mathrm{E}\) 相伴（I, 第 105 页, 引理）。

\subsection{相伴覆叠}

设 B 为拓扑空间，G 为离散拓扑群，E 为 B 的以 G 为群的主覆叠。设 F 为一个 G-集；若给 F 赋予离散拓扑，则群 G 连续地作用于 F。此时 B-空间 \(E \times^{G} F\) 是一个覆叠。它称为与主覆叠 E 相伴的、以 F 为纤维型的 B 的覆叠。

定义 5. —— 设 B 为拓扑空间，G 为离散拓扑群，E 为 B 的以 G 为群的主覆叠。称 B 的覆叠 X 可伴随于主覆叠 E，如果存在一个 G-集 F 以及一个从 E × \(^{G}\) F 到 X 上的 B-同构。

设 B 为拓扑空间，G 为离散拓扑群，E 为 B 的以 G 为群的主覆叠。

对 B 的每个覆叠 X，群 G 从左边作用于 \(\mathcal{C}_{\mathrm{B}}(\mathrm{E};\mathrm{X})\) ，其法则由 \((g\cdot h)(x)=h(x\cdot g)\) 定义，其中 \(x\in E\) ，\(g\in G\) ，\(h\in\mathcal{C}_{\mathrm{B}}(\mathrm{E};\mathrm{X})\) 。

对每个赋予离散拓扑的 G-集 F，由下式定义映射 \(\alpha_{\mathrm{F}}\colon \mathrm{F}\to \mathcal{C}_{\mathrm{B}}(\mathrm{E};\mathrm{E}\times^{\mathrm{G}}\mathrm{F})\) ：

\[
\alpha_ {\mathrm{F}} (y) (x) = \pi (x, y) \quad \text {for } y \in \mathrm{F} \text { and } x \in \mathrm{E},\tag{2}
\]

其中 \(\pi\colon\mathrm{E}\times\mathrm{F}\to\mathrm{E}\times^{\mathrm{G}}\mathrm{F}\) 是典范满射。映射 \(\alpha_{\mathrm{F}}\) 与 G 在 F 上以及 \(\mathcal{C}_{\mathrm{B}}(\mathrm{E};\mathrm{E}\times^{\mathrm{G}}\mathrm{F})\) 上的作用相容。

对 B 的每个覆叠 X，存在唯一的映射 \(\beta_{X}\) （从 \(E \times^{G} \mathcal{C}_{B}(E; X)\) 到 X 中），使得：

\[
\beta_ {\mathrm{X}} (\pi (x, h)) = h (x) \quad \text {for } h \in \mathcal {C} _ {\mathrm{B}} (\mathrm{E}; \mathrm{X}) \text { and } x \in \mathrm{E}.\tag{3}
\]

事实上，有 \((g^{-1}\cdot h)(x\cdot g)=h(x)\) （对 \(x\in E\) ，\(g\in G\) 与 \(h\in\mathcal{C}_{\mathrm{B}}(\mathrm{E};\mathrm{X})\) ），这由 G 在 \(\mathcal{C}_{\mathrm{B}}(\mathrm{E};\mathrm{X})\) 上作用的定义得出。当给集合 \(\mathcal{C}_{\mathrm{B}}(\mathrm{E};\mathrm{X})\) 赋予离散拓扑时，映射 \(\beta_{X}\) 是覆叠的 B-态射。

当 \(\mathrm{X} = \mathrm{E}\times^{\mathrm{G}}\mathrm{F}\) 时，我们有

\[
\beta_ {\mathrm{X}} \left(\pi \left(x, \alpha_ {\mathrm{F}} (y)\right)\right) = \pi (x, y) \quad \text {   for   } x \in \mathrm{X} \text {   and   } y \in \mathrm{F}.\tag{4}
\]

命题 8. —— 设 B 为连通非空拓扑空间。设 G 为离散群，(E, p) 为 B 的以 G 为群的主覆叠。假设 E 连通。沿用上述记号，我们有：

\begin{enumerate}
\def\labelenumi{\alph{enumi})}
\item
  对每个赋予离散拓扑的 G-集 F，映射 \(\alpha_{F}\) 是从 G-集 F 到 G-集 \(\mathcal{C}_{\mathrm{B}}(\mathrm{E};\mathrm{E}\times^{\mathrm{G}}\mathrm{F})\) 上的同构。
\item
  设 X 为 B 的覆叠。B-态射 \(\beta_{X}\) 是从 \(E \times^{G} \mathcal{C}_{B}(E; X)\) 到 X 上的同构，当且仅当 E 的覆叠 \((E \times_{B} X, pr_{1})\) 可平凡化。
\item
  设 \(y\) 与 \(y'\) 为 F 中满足 \(\alpha_{\mathrm{F}}(y) = \alpha_{\mathrm{F}}(y')\) 的点。空间 E 非空；在其中选取一点 \(x\)。我们有 \(\pi(x, y) = \pi(x, y')\) （在 \(\mathrm{E} \times^{\mathrm{G}} \mathrm{F}\) 中）。于是存在 \(g \in \mathrm{G}\) 使得 \(x \cdot g = x\) 且 \(g^{-1} \cdot y = y'\)。前一等式蕴含 \(g\) 是 G 的单位元 \(e\) ，从而 \(y = y'\)。这样，映射 \(\alpha_{\mathrm{F}}\) 是单射。另一方面，设 \(h \in \mathcal{C}_{\mathrm{B}}(\mathrm{E}; \mathrm{E} \times^{\mathrm{G}} \mathrm{F})\) ，并设 \(x\) 为 E 的一点。设 \(x' \in \mathrm{E}\) 与 \(y' \in \mathrm{F}\) 满足 \(h(x) = \pi(x', y')\)。特别地，\(p(x) = p(x')\)；于是存在 G 的元素 \(g\) 使得 \(x' = x \cdot g\) ，并且也有 \(h(x) = \pi(x, y)\) ，其中我们置 \(y = g \cdot y'\)。B-态射 \(h\) 与 \(\alpha_{\mathrm{F}}(y)\) 在 E 的点 \(x\) 处取值相同；由于空间 E 连通，二者相等（I, 第 34 页, 命题 11 的推论 1），这就证明了映射 \(\alpha_{\mathrm{F}}\) 是满射。
\item
  由 I, 第 105 页命题 7 b) 应用于 F = \(\mathcal{C}_{\mathrm{B}}(\mathrm{E};\mathrm{X})\) ，E 上的覆叠 \(p^{*}(\mathrm{E}\times^{\mathrm{G}}\mathcal{C}_{\mathrm{B}}(\mathrm{E};\mathrm{X}))\) 同构于平凡覆叠 \(\mathrm{E}\times \mathcal{C}_{\mathrm{B}}(\mathrm{E};\mathrm{X})\) 。若 \(\beta_{\mathrm{X}}\) 是同构，则 E 上的覆叠 \(p^{*}(\mathrm{X})\) 因之可平凡化。反之，设覆叠 \(p^{*}(\mathrm{X})\) 可平凡化，我们来证明 B-态射 \(\beta_{\mathrm{X}}\) 是双射；由此将得出 \(\beta_{\mathrm{X}}\) 是 B-同构（I, 第 30 页, 命题 6 的推论 2）。
\end{enumerate}

映射 \(\beta_{\mathrm{X}}\) 是由映射 \(\gamma\colon\mathrm{E}\times\mathcal{C}_{\mathrm{B}}(\mathrm{E};\mathrm{X})\to\mathrm{X}\) （由 \(\gamma(x,h)=h(x)\) 定义）过渡到商而得的。我们来证明映射 \(\gamma\) 是满射。设 \(x\) 为 X 的一点。B-空间 E 的投影是满射，因为它是一个主覆叠；于是存在 E 的点 \(y\) 使得 \((y,x)\in\mathrm{E}\times_{\mathrm{B}}\mathrm{X}\) 。于是存在连续截面 \(s\) （可平凡化覆叠 \(\mathrm{pr}_{1}:\mathrm{E}\times_{\mathrm{B}}\mathrm{X}\to\mathrm{E}\) 的截面），使得 \(s(y)=(y,x)\) 。映射 \(h=\mathrm{pr}_{2}\circ s:\mathrm{E}\to\mathrm{X}\) 是 B-态射，且 \(h(y)=x\) ，由此得映射 \(\gamma\) 的满射性，从而得映射 \(\beta_{\mathrm{X}}\) 的满射性。

最后证明 \(\beta_{\mathrm{X}}\) 是单射。设 \((x,h)\) 与 \((x^{\prime},h^{\prime})\) 为 \(\mathrm{E}\times \mathcal{C}_{\mathrm{B}}(\mathrm{E};\mathrm{X})\) 中满足 \(h(x) = h'(x')\) 的元素。注意 \(x\) 与 \(x'\) 在 B 中有同一投影；于是存在 G 的元素 \(g\) 使得 \(x' = x\cdot g\) 。于是有 \(h(x) = h'(x\cdot g) = (g\cdot h')(x)\) 。由于空间 E 连通，有 \(h = g\cdot h'\) （I, 第 34 页, 命题 11 的推论 1）。于是 \((x,h)\) 与 \((x',h')\) 在 \(\mathrm{E}\times^{\mathrm{G}}\mathcal{C}_{\mathrm{B}}(\mathrm{E};\mathrm{X})\) 中有同一类，这表明映射 \(\beta_{\mathrm{X}}\) 是单射，证明完毕。

推论 1. —— 设 (E,p) 为 B 的以 G 为群的主覆叠；假设 E 连通且非空。B 的覆叠 X 可伴随于 E，当且仅当覆叠 \(p^{*}(X)\) 可平凡化。

若覆叠 \(p^{*}(\mathrm{X})\) 可平凡化，则由命题 8 b)，覆叠 X 可伴随于 E。此时，映射 \(\beta_{X}\) 把覆叠 X 等同于与 E 相伴的、以 \(\mathcal{C}_{\mathrm{B}}(\mathrm{E};\mathrm{X})\) 为纤维型的覆叠。反之，设 F 为赋予离散拓扑的 G-集，并设 \(X = E \times^{G} F\) 。此时 \(\alpha_{F}\) 是同构（同前引, a)），且公式 (4) 蕴含 \(\beta_{X}\) 是同构。因此覆叠 \(p^{*}(\mathrm{X})\) 可平凡化（命题 8, b)），推论得证。

推论 2. —— 设 B 为连通且局部连通的拓扑空间，(E,p) 为 B 的以 G 为群的主覆叠。设 E₀ 为 E 的一个连通分支，G₀ 为 G 中稳定 E₀ 的子群。B-空间 (E₀, p \textbar{} E₀) 是以 G₀ 为群的主覆叠（I, 第 103 页, 命题 6）。

每个可伴随于主覆叠 E 的 B 的覆叠 X 都可伴随于主覆叠 \(E_{0}\) 。特别地，覆叠 E 与主覆叠 \(E_{0}\) 相伴。

事实上，若覆叠 X 可伴随于 E，则覆叠 \(p^{*}(\mathrm{X})\) 可平凡化，从而覆叠 \(p_{0}^{*}(\mathrm{X})\) （由 \(p^{*}(\mathrm{X})\) 限制在 \(E_{0}\) 上所得）也可平凡化。

更确切地说，注意映射 \((x,g)\mapsto x\cdot g\) （从 \(E_{0}\times G\) 到 E 中）通过过渡到商，诱导出从 \(E_{0}\times^{G_{0}}G\) 到 E 上的主覆叠的 B-同构。

命题 9. —— 设 B 为拓扑空间，E 为 B 的连通非空、以 G 为群的主覆叠。设 F 为非空的、赋予离散拓扑的 G-集。为使空间 E × \(^{G}\) F 连通，充要条件是 G 传递地作用于 F。

设 U 为 \(E \times^{G} F\) 的既开又闭的子集。若以 \(\pi: E \times F \to E \times^{G} F\) 表示典范满射，则 \(\pi^{-1}(U)\) 是 \(E \times F\) 的在 G 下稳定的既开又闭子集。由于 E 连通，存在在 G 下稳定的子集 \(F' \subset F\)，使得 \(\pi^{-1}(U) = E \times F'\) ，从而 \(U = \pi(E \times F')\)。设 \(F'\) 与 \(F''\) 为 F 的在 G 下稳定的子集，满足 \(\pi(E \times F') = \pi(E \times F'')\)。由于 E 非空，有 \(F' = F''\)。映射 \(F' \mapsto \pi(E \times F')\) 是从 F 的 G 稳定子集的集合到 \(E \times^{G} F\) 的既开又闭子集的集合上的双射。命题得证。

命题 10. —— 设 B 为拓扑空间，(E,p) 为 B 的以 G 为群的主覆叠，X 为 B 的覆叠。假设 E 与 X 都连通且非空。下列性质等价：

\begin{enumerate}
\def\labelenumi{(\roman{enumi})}
\item
  覆叠 X 与主覆叠 E 相伴；
\item
  存在 G 的子群 H，使得 X B-同构于 E/H；
\item
  存在满的 B-态射 \(h: \mathrm{E} \to \mathrm{X}\) ；
\item
  对每个点 \((y, x)\) （属于 \(E \times_{B} X\) ），存在 B-态射 \(r: E \to X\) 使得 \(r(y) = x\) 。
\end{enumerate}

设这些条件成立，且 H 为使 X B-同构于 E/H 的 G 的子群。覆叠 X 是伽罗瓦的，当且仅当子群 H 在 G 中正规。

(i)⇒(ii)：设 F 为使覆叠 \(E \times^{G} F\) B-同构于 X 的离散 G-集。集合 F 非空，且空间 \(E \times^{G} F\) 连通，故群 G 传递地作用于 F（命题 9）。于是空间 \(E \times^{G} F\) B-同构于 E/H，其中 H 是 G 中稳定 F 的某个点的子群（I, 第 106 页, 例 4）。

(ii)⇒(iii)：事实上，从 E 到 E/H 上的典范满射是满的 B-态射。

(iii)⇒(iv)：设 \((y, x)\) 为 \(E \times_{B} X\) 的一点。由于映射 h 是满射，存在 E 的点 \(y'\) 使得 \(h(y') = x\)。点 y 与 \(y'\) 在 B 中有同一投影。由于覆叠 E 是以 G 为群的主覆叠，存在 \(g \in G\) 使得 \(y \cdot g = y'\)。映射 \(r: E \to X\) 由 \(z \mapsto h(z \cdot g)\) 所定义，它是 B-态射，且 \(r(y) = x\)。

(iv)⇒(i)：由于 X 非空且 E 到 B 的投影是满射，空间 \(E \times_{B} X\) 非空。于是存在从 E 到 X 的 B-态射 r。映射 ( \(Id_{E}, r\) ): \(E \to E \times_{B} X\) 是 E 的覆叠 \(p^{*}(X)\) 的一个截面。由于空间 E 连通，由 I, 第 69 页命题 1 的推论 2，覆叠 \(p^{*}(X)\) 可平凡化，这就证明覆叠 X 可伴随于伽罗瓦覆叠 p（命题 8）。

现在设这些条件成立。以 X 表示 B-空间 E/H，q 表示其投影，\(h: E \rightarrow E/H\) 表示典范映射。

若群 H 在 G 中正规，则 X 是以 G/H 为群的主覆叠（I, 第 106 页, 例 4）。由于 X 与 B 连通且非空，X 是 B 的伽罗瓦覆叠（I, 第 102 页, 定理 2）。反之，设 E/H 是 B 的伽罗瓦覆叠，并令 K = Aut \(_{B}\) (E/H) \(^{\circ}\) 。我们定义映射 \(\alpha: \mathrm{E} \times \mathrm{G} \to \mathrm{K}\) ，它把 \((t, g) \in \mathrm{E} \times \mathrm{G}\) 对应到 K 的唯一元素 \(k\) ，后者满足 \(h(t \cdot g) = h(t) \cdot k\) 。它是连续的，因为它是把连续映射 \((t, g) \mapsto (h(t), h(t \cdot g))\) （从 E \(\times\) G 到 X \(\times_{\mathrm{B}}\) X 中）与典范平凡化（I, 第 95 页, 注 1）X \(\times_{\mathrm{B}}\) X \(\to\) X \(\times\) K （即 \(q^{*}(\mathrm{X})\) 的典范平凡化）复合而得。由于 E 连通且 K 是离散群，连续映射 \(t \mapsto \alpha(t, g)\) 对每个 \(g \in \mathrm{G}\) 是常值的；我们将其值记作 \(\alpha(g)\)。若 \(t \in \mathrm{E}\) ，\(g \in \mathrm{G}\) ，且 \(g' \in \mathrm{H}\) ，则有

\[
h (t) = h \left(t \cdot g ^ {- 1} g\right) = h \left(t \cdot g ^ {- 1}\right) \cdot \alpha (g) = h \left(t \cdot g ^ {- 1} \cdot g ^ {\prime}\right) \cdot \alpha (g) = h \left(t \cdot \left(g ^ {- 1} g ^ {\prime} g\right)\right).
\]

由此得出 \(g^{-1}g'g\) 属于 H，因而 H 在 G 中正规。

推论. —— 设 E 为 B 的伽罗瓦覆叠，X 为 B 的连通非空覆叠。假设 X 是有限覆叠，或者空间 B 局部连通。若存在 B-态射 h: E → X，则覆叠 X 与主覆叠 E 相伴。

在两种情形下，B-态射 h 都是覆叠（I, 第 77 页, 定理 1 的推论 3，以及 I, 第 78 页, 命题 5），而连通空间的非空覆叠是满射（I, 第 68 页）。

定理 3. —— 设 B 为非空、连通且局部连通的拓扑空间。B 的每个覆叠都可相伴于某个伽罗瓦覆叠；B 的每个有限覆叠都可相伴于某个有限伽罗瓦覆叠。

设 X 为 B 的覆叠，以 q 表示其投影。由于空间 B 连通且非空，覆叠 X 具有次数；以 F 表示此次数，并给集合 F 赋予离散拓扑。由于空间 B 局部连通，层 \(\mathcal{I} = \underline{\mathcal{I}som}_{\mathrm{B}}(\mathrm{B} \times \mathrm{F}; \mathrm{X})\) 是局部常值的，且在 B 的每点 b 处，其茎 \(I_{b}\) 典范地同构于集合 \(\mathcal{B}(\mathrm{F}; \mathrm{X}_{b})\) ，即从 F 到纤维 \(X_{b}\) 上的双射组成的集合（I, 第 89 页, 命题 12）。设 \(E = E_{\mathcal{I}}\) 为相伴的覆叠（I, 第 86 页, 推论）。

设 G 为 F 的置换群。对每个 \(g \in G\)，我们如下定义一个态射 \(\gamma'(g)\) （层 \(\mathcal{I}\) 到自身的态射）：对 B 的每个开集 U，映射 \(\gamma'(g)_{\mathrm{U}}\) 把 U-同构 \(\varphi: \mathrm{U} \times \mathrm{F} \to q^{-1}(\mathrm{U})\) 对应到由 \((b, f) \mapsto \varphi(b, g(f))\) （其中 \(b \in U\)，\(f \in F\)）定义的 U-同构。我们有 \(\gamma'(\mathrm{Id}_{\mathrm{F}}) = \mathrm{Id}_{\mathcal{I}}\) 以及 \(\gamma'(g \circ g') = \gamma'(g') \circ \gamma'(g)\) （其中 \(g, g' \in G\)），故对每个 \(g \in G\)，\(\gamma'(g)\) 是层 \(\mathcal{I}\) 的自同构。B-态射 \(\gamma(g): E \to E\) （与 \(\gamma'(g)\) 相伴）是自同构（I, 第 50 页）。若 \(x = [\mathrm{U}, \varphi, b]\) 是 E 的元素，其中 U 是 B 的开子集，b 是 U 的点，\(\varphi\) 是从 U×F 到 \(q^{-1}(\mathrm{U})\) 上的 U-同构，则有 \(\gamma(g)(x) = [\mathrm{U}, \psi, b]\) ，其中 \(\psi\) 由 \(\psi(a, f) = \varphi(a, g(f))\) 定义（\(a \in U\)，\(f \in F\)）。若 g 与 \(g'\) 是 G 的元素，则有 \(\gamma(g \circ g') = \gamma(g') \circ \gamma(g)\)。我们有 \(\gamma(\mathrm{Id}_{\mathrm{F}}) = \mathrm{Id}_{\mathrm{E}}\)。于是，通过置 \(x \cdot g = \gamma(g)(x)\) （其中 \(x \in E\)，\(g \in G\)），我们定义了 G 在 E 上的连续右作用。群 G 单传递地作用于 E 的每条纤维，故赋予此作用的覆叠 E 是以 G 为群的主覆叠（I, 第 97 页, 命题 4）。设 h 为从 E × F 到 X 中、由 \(h([\mathrm{U}, \varphi, b], f) = \varphi(b, f)\) 定义的映射，其中 U 取遍 B 的开子集，b 取遍 U 的点，\(\varphi\) 取遍从 U × F 到 \(q^{-1}(\mathrm{U})\) 上的 U-同构。由 E 上拓扑的定义，它是连续的；它是一个 B-态射。对 G 的每个元素与 E × F 的每个点 \((x, f)\)，有 \(h(x, g(f)) = h(x \cdot g, f)\)。因此映射 h 通过过渡到商，定义一个 B-态射 \(h': E \times^{G} F \to X\)。对 B 的每个点 \(b\) ，映射 \(h_b\colon \mathrm{E}_b\times \mathrm{F}\to \mathrm{X}_b\) 等同于映射 \((\varphi, f)\mapsto \varphi(f)\) （从 \(\mathcal{B}(\mathrm{F};\mathrm{X}_b)\times \mathrm{F}\) 到 \(\mathrm{X}_b\) 中），故映射 \(h_b^\prime\) 是双射。因此，\(h^\prime\) 是从 \(\mathrm{E}\times^{\mathrm{G}}\mathrm{F}\) 到 X 上的 B-同构（I, 第 30 页, 命题 6 的推论 2）。

由于空间 B 连通、局部连通且非空，可伴随于主覆叠 E 的覆叠 X 可相伴于某个伽罗瓦覆叠（I, 第 110 页, 推论 2）。若覆叠 X 有限，则覆叠 E 亦然，故 X 可相伴于某个有限伽罗瓦覆叠（同前引）。

\subsection{由上闭链定义的主纤维丛}

定义 6. —— 设 B 为拓扑空间，G 为拓扑群，\(\mathcal{U} = (\mathrm{U}_{i})_{i\in \mathrm{I}}\) 为 B 的开覆盖。B 上取值于 G、从属于 \(\mathcal{U}\) 的连续 1-上闭链，是指对 I 的元素的每个偶 \((i,j)\) ，给定一个连续映射 \(g_{i,j}\) （从 \(\mathrm{U}_i\cap \mathrm{U}_j\) 到 G 中），使得对每个三元组 \((i,j,k)\in \mathrm{I}\times \mathrm{I}\times \mathrm{I}\) 以及 \(\mathrm{U}_i\cap \mathrm{U}_j\cap \mathrm{U}_k\) 的每个点 b，有

\[
g _ {i, k} (b) = g _ {i, j} (b) g _ {j, k} (b).
\]

我们以 \(Z_{\mathrm{cont}}^{1}(\mathrm{B},\mathcal{U},\mathrm{G})\) 表示 B 上取值于 G、从属于覆盖 U 的连续 1-上闭链组成的集合。

在本分册中，我们把连续 1-上闭链简称为上闭链。

设 \((\mathrm{E},p)\) 为 B 上的主 G-丛，\(\mathcal{U} = (\mathrm{U}_{i})_{i\in\mathrm{I}}\) 为由开集 \(U_{i}\) 组成的 B 的覆盖。若对每个 \(i\in I\)，E 在 \(U_{i}\) 上都可平凡化，则称 E 在 U 上可平凡化。此时我们把 E 的 U-平凡化称为族 \((f_{i})_{i\in\mathrm{I}}\) ，其中 \(f_{i}\colon p^{-1}(\mathrm{U}_{i})\to\mathrm{U}_{i}\times\mathrm{G}\) 是 E 在 \(U_{i}\) 上的一个平凡化。

设 \(\mathcal{U} = (\mathrm{U}_i)_{i\in \mathrm{I}}\) 为由开集组成的 B 的覆盖，\((\mathrm{E},p)\) 为 B 上赋予了 \(\mathcal{U}\)-平凡化 \((f_i)_{i\in \mathrm{I}}\) 的主 G-丛。对每个偶 \((i,j)\in \mathrm{I}\times \mathrm{I}\) ，以 \(g_{ij}\) 表示从 \(\mathrm{U}_i\cap \mathrm{U}_j\) 到 G 中由

\[
(b, g _ {i j} (b)) = f _ {i} \circ f _ {j} ^ {- 1} (b, e),
\]

所定义的映射，其中 e 表示 G 的单位元。它是连续的。

由于 \(f_{i}\) 与 \(f_{j}\) 与 G 的作用相容，我们有

\[
(f _ {i} \circ f _ {j} ^ {- 1}) (b, g) = (b, g _ {i j} (b) g)
\]

其中 \(b \in \mathrm{U}_i \cap \mathrm{U}_j\) ，\(g \in \mathrm{G}\) 。若 \(b\) 是 \(\mathrm{U}_i \cap \mathrm{U}_j \cap \mathrm{U}_k\) 的点（其中 \(i, j, k \in \mathrm{I}\)），则有

\[
(f _ {i} \circ f _ {k} ^ {- 1}) (b, e) = (b, g _ {i k} (b))
\]

以及

\[
\begin{array}{c} (f _ {i} \circ f _ {k} ^ {- 1}) (b, e) = (f _ {i} \circ f _ {j} ^ {- 1}) \circ (f _ {j} \circ f _ {k} ^ {- 1}) (b, e) \\ = (f _ {i} \circ f _ {j} ^ {- 1}) (b, g _ {j k} (b)) = (b, g _ {i j} (b) g _ {j k} (b)). \end{array}
\]

由此得出

\[
g _ {i k} (b) = g _ {i j} (b) g _ {j k} (b),
\]

故族 \((g_{ij})\) 是 B 上取值于 G、从属于覆盖 U 的上闭链。它称为由平凡化族 \((f_{i})_{i\in\mathrm{I}}\) 定义的上闭链。

设 B、G 与 U 如定义 6 所述，\((g_{ij}) \in Z_{\mathrm{cont}}^{1}(\mathrm{B}, \mathcal{U}, \mathrm{G})\) 为一个上闭链。对每个偶 \((i, j) \in \mathrm{I} \times \mathrm{I}\) ，映射 \(\gamma_{ij} \colon (\mathrm{U}_{i} \cap \mathrm{U}_{j}) \times \mathrm{G} \to (\mathrm{U}_{i} \cap \mathrm{U}_{j}) \times \mathrm{G}\) 由

\[
\gamma_ {i j} (b, g) = \left(b, g _ {i j} (b) g\right) \quad \text {   for   } b \in \mathrm{U} _ {i} \cap \mathrm{U} _ {j} \text {   and   } g \in \mathrm{G}\tag{5}
\]

所定义，它是基 \(U_{i} \cap U_{j}\) 上主丛的同构。对每个三元组 \((i,j,k) \in \mathrm{I} \times \mathrm{I} \times \mathrm{I}\) 以及每个偶 \((b,g) \in (\mathrm{U}_{i} \cap \mathrm{U}_{j} \cap \mathrm{U}_{k}) \times \mathrm{G}\) ，我们有：

\[
\gamma_ {i k} (b, g) = \gamma_ {i j} \circ \gamma_ {j k} (b, g).
\]

设 F 为把空间 \(U_{i} \times G\) 沿 \((\mathrm{U}_{i} \cap \mathrm{U}_{j}) \times \mathrm{G}\) 借助双射 \(\gamma_{ij}\) 粘合起来所得的拓扑空间（TG, I, 第 16 页）。对每个 \(i \in I\)，\(U_{i} \times G\) 在 F 中的像是 F 的开集（TG, I, 第 17 页, 命题 9）。典范投影 \(p_{i}: U_{i} \times G \to U_{i}\) 粘合成一个从 F 到 B 的连续映射 p。由于诸映射 \(\gamma_{ij}\) 与空间 \(U_{i} \times G\) 中的 G 右作用相容，由此导出 G 在 F 上的连续右作用，它使 F 成为以 B 为基、以 G 为群、且在每个 \(U_{i}\) 上带有平凡化的主纤维丛。我们称 F 为由上闭链 \((g_{ij})\) 定义的主纤维丛。

定义 7. —— 设 B 为拓扑空间，G 为拓扑群，\(\mathcal{U} = (\mathrm{U}_{i})_{i\in \mathrm{I}}\) 为 B 的开覆盖。称两个上闭链 \((g_{ij})\) 与 \((g_{ij}^{\prime})\) （均属于 \(\mathrm{Z}_{\mathrm{cont}}^{1}(\mathrm{B},\mathcal{U},\mathrm{G})\) ）是上同调的，如果存在族 \((h_{i})_{i\in\mathrm{I}}\) （由连续映射 \(h_{i}\colon U_{i}\to G\) 组成），使得

\[
g _ {i j} ^ {\prime} (b) = h _ {i} (b) g _ {i j} (b) h _ {j} (b) ^ {- 1}\tag{6}
\]

对每个偶 \((i,j)\in \mathrm{I}\times \mathrm{I}\) 以及每个 \(b\in \mathrm{U}_i\cap \mathrm{U}_j\) 成立。

关系"\((g_{ij})\) 与 \((g'_{ij})\) 上同调"是集合 \(Z_{\mathrm{cont}}^{1}(\mathrm{B},\mathcal{U},\mathrm{G})\) 上的一个等价关系。我们以 \(H_{\mathrm{cont}}^{1}(\mathrm{B},\mathcal{U},\mathrm{G})\) 表示 \(Z_{\mathrm{cont}}^{1}(\mathrm{B},\mathcal{U},\mathrm{G})\) 关于此等价关系的商集。

命题 11. —— 设 B 为拓扑空间，G 为拓扑群，\(\mathcal{U} = (\mathrm{U}_{i})_{i\in \mathrm{I}}\) 为 B 的开覆盖。

\begin{enumerate}
\def\labelenumi{\alph{enumi})}
\item
  B 上每个在 U 上可平凡化的主 G-丛都同构于由 \(Z_{\mathrm{cont}}^{1}(\mathrm{B}, \mathcal{U}, \mathrm{G})\) 中某个上闭链定义的主丛。
\item
  设 (E, p) 与 (E', p') 为 B 上在 U 上可平凡化的主 G-丛。设 \((f_i)_{i \in I}\) （相应地 \((f_i')_{i \in I}\) ）为 (E, p) （相应地 \((E', p')\) ）的适合 U 的平凡化，并以 \((g_{ij})_{(i,j) \in I \times I}\) （相应地 \((g_{i,j}')_{(i,j) \in I \times I}\) ）表示由此平凡化定义的上闭链。此时，主纤维丛 (E, p) 与 (E', p') 同构，当且仅当这些上闭链上同调。
\end{enumerate}

我们来证明 \(b\) 。设 \(\varphi: \mathrm{E} \to \mathrm{E}'\) 为以 B 为基、以 G 为群的主丛的同构。对每个 \(i \in \mathrm{I}\) ，设 \(h_i\) 为从 \(\mathrm{U}_i\) 到 G 中的连续映射，由

\[
(b, h _ {i} (b)) = f _ {i} ^ {\prime} \circ \varphi \circ f _ {i} ^ {- 1} (b, e).
\]

所定义。由于对每个 \(i \in I\) ，\(f_{i}\) 与 \(f_{i}^{\prime}\) 都与 G 的作用相容，故对每个 \(b \in U_{i}\) 与每个 \(g \in G\) ，有

\[
(f _ {i} ^ {\prime} \circ \varphi \circ f _ {i} ^ {- 1}) (b, g) = (b, h _ {i} (b) g)
\]

以及

\[
(f _ {i} \circ \varphi^ {- 1} \circ (f _ {i} ^ {\prime}) ^ {- 1}) (b, g) = (b, h _ {i} (b) ^ {- 1} g).
\]

因此，对每个偶 \((i,j)\in \mathrm{I}\times \mathrm{I}\) 以及每个点 \(b\) （属于 \(\mathrm{U}_i\cap \mathrm{U}_j\) ），

\[
\begin{array}{c} f _ {i} ^ {\prime} \circ (f _ {j} ^ {\prime}) ^ {- 1} (b, e) = (f _ {i} ^ {\prime} \circ \varphi \circ f _ {i} ^ {- 1}) \circ (f _ {i} \circ f _ {j} ^ {- 1}) \circ (f _ {j} \circ \varphi^ {- 1} \circ (f _ {j} ^ {\prime}) ^ {- 1}) (b, e) \\ = (b, h _ {i} (b) g _ {i j} (b) h _ {j} (b) ^ {- 1}) \end{array}
\]

从而有

\[
g _ {i j} ^ {\prime} (b) = h _ {i} (b) g _ {i j} (b) h _ {j} (b) ^ {- 1}.
\]

这就证明上闭链 \((g_{ij})\) 与 \((g'_{ij})\) 上同调。

反之，设这些上闭链上同调，并设 \((h_{i})_{i\in\mathrm{I}}\) 为由连续映射 \(h_{i}\colon U_{i}\to G\) 组成的族，它满足 \(g_{ij}^{\prime}(b)=h_{i}(b)g_{ij}(b)h_{j}(b)^{-1}\) （其中 \(i,j\in I\)，\(b\in U_{i}\cap U_{j}\)）。对 \(i\in I\) ，设 \(\varphi_{i}\colon p^{-1}(U_{i})\to(p')^{-1}(U_{i})\) 为由 \(f_{i}^{\prime}\circ\varphi_{i}\circ f_{i}^{-1}(b,g)=(b,h_{i}(b)g)\) （其中 \(b\in U_{i}\)，\(g\in G\)）定义的映射。这是以 \(U_{i}\) 为基、以 G 为群的主纤维丛的同构。对 \((i,j)\in I\times I\) ，分别以 \(\gamma_{ij}\) 与 \(\gamma_{ij}^{\prime}\) 表示如上依附于上闭链 \((g_{ij})\) 与 \((g_{ij}^{\prime})\) 的粘合同胚（I, 第 115 页, 公式 (5)），使得 \(f_{i}(b,g)=\gamma_{ij}\circ f_{j}(b,g)\) 与 \(f_{i}^{\prime}(b,g)=\gamma_{ij}^{\prime}\circ f_{j}^{\prime}(b,g)\) 对一切 \((b,g)\in(U_{i}\cap U_{j})\times G\) 成立。于是，对 \((i,j)\in I\times I\) 与 \((b,g)\in(U_{i}\cap U_{j})\times G\) ，有关系

\[
\begin{array}{r l} & f _ {i} ^ {\prime} \circ \varphi_ {j} \circ f _ {i} ^ {- 1} (b, g) = \gamma_ {i j} ^ {\prime} \circ (f _ {j} ^ {\prime} \circ \varphi_ {j} \circ f _ {j} ^ {- 1}) (b, g _ {i j} (b) ^ {- 1} g) \\ & \qquad = \gamma_ {i j} ^ {\prime} (b, h _ {j} (b) g _ {i j} (b) ^ {- 1} g) \\ & \qquad = (b, g _ {i j} ^ {\prime} (b) h _ {j} (b) g _ {i j} (b) ^ {- 1} g) \\ & \qquad = (b, h _ {i} (b) g) = f _ {i} ^ {\prime} \circ \varphi_ {i} \circ f _ {i} ^ {- 1} (b, g). \end{array}
\]

这表明 \(\varphi_{i}\) 与 \(\varphi_{j}\) 在 \(p^{-1}(\mathrm{U}_{i}\cap\mathrm{U}_{j})\) 上相合。于是诸态射 \(\varphi_{i}\) 粘合成一个从 E 到 E' 的主丛的 B-态射。断言 \(b\) 由"任何以 B 为基、以 G 为群的主丛态射都是同构"得出（I, 第 93 页, 命题 1）。

现在来证明 \(a\) 。设 \((\mathrm{E}, p)\) 为以 G 为结构群的主纤维丛，且对每个 \(i \in \mathrm{I}\) ，带有平凡化 \(f_i: p^{-1}(\mathrm{U}_i) \to \mathrm{U}_i \times \mathrm{G}\) （在 \(\mathrm{U}_i\) 上）。设 \((g_{ij})\) 为由此族定义的上闭链，F 为由上闭链 \((g_{ij})\) 定义的主纤维丛。按构造，主纤维丛 F 带有 \(\mathcal{U}\) 上的一个平凡化；该平凡化定义出上闭链 \((g_{ij})\)。由断言 \(b\) ，主纤维丛 \((\mathrm{E}, p)\) 同构于 F。

设 B 为拓扑空间，G 为拓扑群。设 \((E,p)\) 为以 B 为基、以 G 为群的主纤维丛。设 s 为满射 p 的一个截面（E, II, 第 19 页, 命题 8）。映射 \(f\colon B\times G\to E\) 由 \(f(b,g)=s(b)\cdot g\) 所定义，它是与 G 的作用相容的双射。给 \(B\times G\) 赋予经结构转移得到的拓扑；此时 \((B\times G,\mathrm{pr}_{1})\) 是以 B 为基、以 G 为群且同构于 E 的主纤维丛。因此，存在一个以 B 为基、以 G 为群的主丛的集合 T，使得每个以 B 为基、以 G 为群的主纤维丛都同构于 T 的某个元素。以 P(B, G) 表示以 B 为基、以 G 为群的主丛的同构类组成的集合（E, II, 第 47 页）。

设 \(\mathcal{U} = (\mathrm{U}_i)_{i\in \mathrm{I}}\) 为 B 的开覆盖。以 \(\mathrm{P(B,\mathcal{U},G)}\) 表示 \(\mathrm{P(B,G)}\) 中由在 \(\mathcal{U}\) 上可平凡化的主丛的同构类组成的子集。由命题 11，存在映射 \(r_{\mathcal{U}}\) （从 \(\mathrm{H}_{\mathrm{cont}}^{1}(\mathrm{B},\mathcal{U},\mathrm{G})\) 到 \(\mathrm{P(B,\mathcal{U},G)}\) 中），它把上闭链 \((g_{ij})\in Z_{\mathrm{cont}}^{1}(\mathrm{B},\mathcal{U},\mathrm{G})\) 的类对应到由此上闭链定义的主丛的同构类；它是双射（同前引）。逆双射 \(s_{\mathcal{U}}\) 把 \(\mathrm{P(B,\mathcal{U},G)}\) 中某个在 \(\mathcal{U}\) 上可平凡化的主纤维丛 E 的同构类，对应到由 E 在 \(\mathcal{U}\) 上任一平凡化族所定义的上闭链的类。在此对应下，平凡主纤维丛 \(\mathrm{B}\times \mathrm{G}\) 的同构类对应于常值上闭链 \(g_{ij} = e\) 的上同调类，后者也称为平凡上闭链。

依据主纤维丛的定义，集合 P(B, G) 是形如 P(B, U, G) 的集合之并，其中 U 取遍 B 的开覆盖。

设 \(\mathcal{U} = (\mathrm{U}_i)_{i\in \mathrm{I}}\) 为 B 的开覆盖，\(\mathcal{V} = (\mathrm{V}_k)_{k\in \mathrm{K}}\) 为比 \(\mathcal{U}\) 更细的 B 的开覆盖。我们有 \(\mathrm{P(B,\mathcal{U},G)}\subset\) \(\mathrm{P(B,\mathcal{V},G)}\) ；以 \(i_{\mathcal{V}\mathcal{U}}\) 表示由此包含关系定义的典范单射。选取一个映射 \(\varphi \colon \mathrm{K}\to \mathrm{I}\) ，使得 \(\mathrm{V}_k\subset \mathrm{U}_{\varphi (k)}\) 对一切 \(k\in \mathrm{K}\) 成立。给定上闭链 \((g_{ij})\in \mathrm{Z}_{\mathrm{cont}}^1 (\mathrm{B},\mathcal{U},\mathrm{G})\) ，对每个偶 \((k,\ell)\in \mathrm{K}\times \mathrm{K}\) ，置 \(\overline{g}_{k\ell} = g_{\varphi (k)\varphi (\ell)}\mid \mathrm{V}_k\cap \mathrm{V}_\ell\) 。

族 \((\overline{g}_{k\ell})\) 是 B 上取值于 G、从属于 V 的上闭链。若 \((g'_{ij}) \in \mathrm{Z}_{\mathrm{cont}}^{1}(\mathrm{B}, \mathcal{U}, \mathrm{G})\) 是与上闭链 \((g_{ij})\) 上同调的上闭链，则上闭链 \((\overline{g}'_{k\ell})\) （由 \((g'_{ij})\) 得出）与上闭链 \((\overline{g}_{k\ell})\) 上同调。这给出映射

\[
c (\varphi) \colon \mathrm{H} _ {\mathrm{cont}} ^ {1} (\mathrm{B}, \mathcal {U}, \mathrm{G}) \to \mathrm{H} _ {\mathrm{cont}} ^ {1} (\mathrm{B}, \mathcal {V}, \mathrm{G})
\]

它把 \((g_{ij})\) 的类对应到 \((\overline{g}_{k\ell})\) 的类。

设 E 为以 B 为基、以 G 为群的主纤维丛，且对每个 \(i \in I\) ，设 \(f_i: E_{U_i} \to U_i \times G\) 为 \(U_i\)-主纤维丛 \(E_{U_i}\) 的一个平凡化。对每个 \(k \in K\) ，设 \(f_k' : E_{V_k} \to V_k \times G\) 为由 \(f_{\varphi(k)}\) 限制到子集而得的平凡化。设 \((g_{ij}) \in Z_{\mathrm{cont}}^1(\mathrm{B}, \mathcal{U}, \mathrm{G})\) 为由族 \((f_i)\) 定义的上闭链；由族 \((f_{k}^{\prime})\) 定义的上闭链正是上面定义的上闭链 \((\overline{g}_{k\ell})\) 。于是，下面的图式是交换的：

\[
\begin{array}{c} \mathrm{H} _ {\mathrm{cont}} ^ {1} (\mathrm{B}, \mathcal {U}, \mathrm{G}) ^ {c (\varphi)} \longrightarrow \mathrm{H} _ {\mathrm{cont}} ^ {1} (\mathrm{B}, \mathcal {V}, \mathrm{G}) \\ \Big \downarrow^ {r _ {\mathcal {U}}} \qquad \qquad \qquad \qquad \Big \downarrow^ {r _ {\mathcal {V}}} \\ \mathrm{P} (\mathrm{B}, \mathcal {U}, \mathrm{G}) ^ {i _ {\mathcal {V} \mathcal {U}}} \longrightarrow \mathrm{P} (\mathrm{B}, \mathcal {V}, \mathrm{G}). \end{array}
\]

由于映射 \(r_{\mathcal{U}}\) 与 \(r_{\mathcal{V}}\) 都是双射，映射 \(c(\varphi)\) 与 \(i_{\mathcal{V}\mathcal{U}}\) 一样是单射，且不依赖于 \(\varphi\) 的选取。此后我们将以 \(c_{\mathcal{V}\mathcal{U}}\) 代替 \(c(\varphi)\) 来记写。

设 \(\mathcal{R}\) 为 \(\mathfrak{P}(\mathfrak{P}(B))\) 中构成 B 的开覆盖的元素组成的集合。对 B 的每个开覆盖 \(\mathcal{U}\) ，存在开覆盖 \(\mathcal{V}\) （属于 \(\mathcal{R}\) ），使得 \(\mathcal{V}\) 比 \(\mathcal{U}\) 既更细又更粗。集合 \(\mathcal{R}\) 对于关系 \(\leqslant\) （由"\(\mathcal{U} \leqslant \mathcal{V}\) 当且仅当 \(\mathcal{V}\) 是比 \(\mathcal{U}\) 更细的覆盖"定义）是有序的且是滤子的。由前所述，我们已定义了归纳系 \((\mathrm{H}_{\mathrm{cont}}^{1}(\mathrm{B},\mathcal{U},\mathrm{G}),c_{\mathcal{V}\mathcal{U}})\) （关于滤子化有序集 \(\mathcal{R}\) ），且族 \((r_{\mathcal{U}})\) 是从 \((\mathrm{H}_{\mathrm{cont}}^{1}(\mathrm{B},\mathcal{U},\mathrm{G}),c_{\mathcal{V}\mathcal{U}})\) 到 \((\mathrm{P}(\mathrm{B},\mathcal{U},\mathrm{G}),i_{\mathcal{V}\mathcal{U}})\) 的双射映射的归纳系。若以 \(\mathrm{H}_{\mathrm{cont}}^{1}(\mathrm{B},\mathrm{G})\) 表示系 \((\mathrm{H}_{\mathrm{cont}}^{1}(\mathrm{B},\mathcal{U},\mathrm{G}),c_{\mathcal{V}\mathcal{U}})\) 的归纳极限，以 \(r\colon \mathrm{H}_{\mathrm{cont}}^{1}(\mathrm{B},\mathrm{G})\to \mathrm{P}(\mathrm{B},\mathrm{G})\) 表示族 \((r_{\mathcal{U}})\) 的归纳极限，则有：

定理 4. —— 映射 \(r\colon \mathrm{H}_{\mathrm{cont}}^{1}(\mathrm{B},\mathrm{G})\to \mathrm{P}(\mathrm{B},\mathrm{G})\) 是双射。

设 \(\mathcal{U} = (\mathrm{U}_i)_{i\in \mathrm{I}}\) 为 B 的开覆盖；以 \(c_{\mathcal{U}}\) 表示典范映射 \(\mathrm{H}_{\mathrm{cont}}^{1}(\mathrm{B},\mathcal{U},\mathrm{G})\to \mathrm{H}_{\mathrm{cont}}^{1}(\mathrm{B},\mathrm{G})\) 。若 \((g_{ij})\) 是 B 上取值于 G、从属于开覆盖 \(\mathcal{U}\) 的上闭链，则元素 \(c_{\mathcal{U}}((g_{ij}))\) （属于 \(\mathrm{H}_{\mathrm{cont}}^{1}(\mathrm{B},\mathrm{G})\) ）称为上闭链 \((g_{ij})\) 的上同调类。

\section{单连通空间}

\subsection{万有覆叠}

定义 1. —— 带基点集是指带有其一个元素的集合 X，该元素称为基点。带有元素 x 的集合 X 有时记作 (X, x)。设 (X, x) 与 (Y, y) 为带基点集；从 (X, x) 到 (Y, y) 的带基点映射是指满足 \(f(x) = y\) 的从 X 到 Y 的映射 f。

带基点拓扑空间、带基点连续映射、带基点拓扑空间的带基点覆叠等概念，都可类似地定义。

若 \((X,x)\) 与 \((Y,y)\) 为带基点拓扑空间，则从 \((X,x)\) 到 \((Y,y)\) 的带基点连续映射组成的集合记作 \(\mathcal{C}((X,x);(Y,y))\)。

人们常不用"设 f 为从 \((X, x)\) 到 \((Y, y)\) 上的带基点映射"这种说法，而使用下述短语："设 \(f: (X, x) \to (Y, y)\) 为带基点映射"。

设 B 为拓扑空间，b 为 B 的一点。

定义 2. —— 称 (B,b) 的带基点覆叠 (E,x) 为万有覆叠，如果对 (B,b) 的每个带基点覆叠 (E',x')，存在唯一的 B 的覆叠态射 \(f: \mathrm{E} \to \mathrm{E}'\) ，使得 \(f(x) = x'\) 。

若 (E, x) 与 (E', x') 都是 (B, b) 的万有覆叠，则从 (E, x) 到 (E', x') 的唯一的 B-态射是 B-空间的同构。

设 E 为 B 的连通覆叠，x 为纤维 \(E_{b}\) 的一点。假设对每个带基点覆叠 \((\mathrm{E}^{\prime}, x^{\prime})\) （\((\mathrm{B}, b)\) 的），都存在 B 的覆叠态射 \(f: E \to E^{\prime}\) 使得 \(f(x) = x^{\prime}\) 。这样的态射 f 于是是唯一的（I, 第 34 页, 命题 11 的推论 1），故 \((\mathrm{E}, x)\) 是 \((\mathrm{B}, b)\) 的万有覆叠。*我们以后将会看到（I, 第 126 页, 命题 3 的推论），特别地，当 E 的一切覆叠都可平凡化时即属此情形。*

命题 1. —— 设 B 为连通且局部连通的拓扑空间，b 为 B 的一点。设 (E, x) 为 (B, b) 的万有覆叠。则 E 是 B 的伽罗瓦覆叠，且 B 的每个覆叠都与 E 相伴。

空间 E 是局部连通的，因为 B 是局部连通的。设 \(E_{0}\) 为 x 在 E 中的连通分支，于是空间 \((\mathrm{E}_{0}, x)\) 是 \((\mathrm{B}, b)\) 的带基点覆叠（I, 第 80 页, 命题 6 的推论 1）。于是存在唯一的覆叠态射 \(f: E \to E_{0}\) 使得 \(f(x) = x\)。若以 i 表示 \(E_{0}\) 到 E 中的典范单射，则映射 \(i \circ f: E \to E\) 是把 x 映到 x 的覆叠态射，映射 \(Id_{E}\) 亦然；由于 \((\mathrm{E}, x)\) 是 \((\mathrm{B}, b)\) 的万有覆叠，我们有 \(i \circ f = Id_{E}\)。这蕴含 i 是满射，从而 \(E_{0} = E\)。于是 E 是连通的。

设 \(y\) 为 \(\mathrm{E}_b\) 的一点，并考察带基点覆叠 \((\mathrm{E},y)\) （\((\mathrm{B},b)\) 的）；依假设存在唯一的覆叠态射 \(f\colon \mathrm{E}\to \mathrm{E}\) 使得 \(f(x) = y\) 。映射 \(s\colon \mathrm{E}\to \mathrm{E}\times_{\mathrm{B}}\mathrm{E}\) 由 \(t\mapsto (t,f(t))\) 所定义，它是映射 \(\mathrm{pr}_1\colon \mathrm{E}\times_{\mathrm{B}}\mathrm{E}\to \mathrm{E}\) 的连续截面。由 I, 第 81 页推论 4，这表明覆叠 \(\mathrm{E}\times_{\mathrm{B}}\mathrm{E}\) （由映射 \(\mathrm{pr}_1\) 给出的 E 上的覆叠）可平凡化。于是覆叠 E 是伽罗瓦的（I, 第 102 页, 定理 2）。再由 I, 第 112 页命题 10 的推论，B 的每个覆叠都与 E 相伴。

推论. —— 设 B 为局部连通拓扑空间。设 b 为 B 的一点，(E, x) 为 (B, b) 的万有覆叠。对 B 的子空间 A，下面两条性质等价：

\begin{enumerate}
\def\labelenumi{(\roman{enumi})}
\item
  覆叠 E 在 A 上可平凡化；
\item
  B 的每个覆叠都在 A 上可平凡化。
\end{enumerate}

性质 (ii) 显然蕴含性质 (i)。其逆由"B 的每个覆叠都与覆叠 E 相伴"这一事实得出（I, 第 105 页, 命题 7）。

\subsection{数值空间中的凸子集}

设 E 为 n 维数值空间\(^{*}\)（或更一般地，R 上的拓扑向量空间）*。对 E 的点的每个偶 \((x,y)\) ，称形如 \(tx + (1-t)y\) （其中 \(t \in [0,1]\) ；相应地 \(t \in ]0,1[\) ）的点组成的集合为以 x 与 y 为端点的线段（相应地，开线段）（参见 EVT, II, 第 7 页）。设 A 为 E 的子集。若对 A 的点的每个偶 \((x,y)\) 以及每个 \(t \in I\) ，点 \(tx + (1-t)y\) 都属于 A，则称集合 A 是凸的。

*凸子集是道路连通的。*

引理 1. —— 设 E 为 n 维数值空间，A 为 E 的凸紧子集，且 0 是它的一个内点。对每个 \(x \in E\) ，以 \(p_{\mathrm{A}}(x)\) 表示在 \(\overline{R}\) 中所取的、由实数 \(t > 0\) （使 \(x \in tA\) ）组成的集合的下确界。

映射 \(p_{A}\) 是有限的、连续的，且满足下列性质：

\begin{enumerate}
\def\labelenumi{(\roman{enumi})}
\item
  对每个 \(x \in E\) （满足 \(x \neq 0\) ），有 \(p_{\mathrm{A}}(x) > 0\) ；
\item
  对每个 \(s \in R_{+}\) 与每个 \(x \in E\) ，有 \(p_{\mathrm{A}}(sx) = sp_{\mathrm{A}}(x)\) 。
\item
  对一切 \(x\) 与 \(y \in \mathrm{E}\) ，有 \(p_{\mathrm{A}}(x + y) \leqslant p_{\mathrm{A}}(x) + p_{\mathrm{A}}(y)\) 。
\item
  为使 E 的点 \(x\) 属于 A，充要条件是 \(p_{\mathrm{A}}(x) \leqslant 1\) 。
\end{enumerate}

对 \(x \in E\) ，以 \(\|x\|\) 表示 x 的欧氏范数（TG, VI, 第 7 页）。由于 A 是紧的，存在实数 \(M > 0\)，使得 A 的每个点 x 满足 \(\|x\| \leqslant M\)。由于 0 是 A 的内点，存在实数 \(m > 0\)，使得 E 中满足 \(\|x\| \leqslant m\) 的每个点 x 都属于 A。于是有关系 \(\|x\|/\mathrm{M} \leqslant p_{\mathrm{A}}(x) \leqslant \|x\|/m\) （对每个 \(x \in E\)）。特别地，\(p_{\mathrm{A}}(0) = 0\) ，且 \(p_{\mathrm{A}}(x) \neq 0\) （若 \(x \neq 0\)）。

断言 (ii) 与 (iv) 由映射 \(p_{A}\) 的定义立得。

设 \(x\) 与 \(y\) 为 E 的点。设 \(x'\) 与 \(y'\) 为 A 中满足 \(x = p_{\mathrm{A}}(x)x'\) 与 \(y = p_{\mathrm{A}}(y)y'\) 的点。若 \(x\) 与 \(y\) 不全为零，则 \(p_{\mathrm{A}}(x) + p_{\mathrm{A}}(y) > 0\) ，且有

\[
\begin{array}{l} x + y = p _ {\mathrm{A}} (x) x ^ {\prime} + p _ {\mathrm{A}} (y) y ^ {\prime} \\ \qquad = (p _ {\mathrm{A}} (x) + p _ {\mathrm{A}} (y)) \left(\frac {p _ {\mathrm{A}} (x)}{p _ {\mathrm{A}} (x) + p _ {\mathrm{A}} (y)} x ^ {\prime} + \frac {p _ {\mathrm{A}} (y)}{p _ {\mathrm{A}} (x) + p _ {\mathrm{A}} (y)} y ^ {\prime}\right). \end{array}
\]

由于 A 是凸的，这表明 \(x + y\) 属于 \((p_{\mathrm{A}}(x) + p_{\mathrm{A}}(y))\mathrm{A}\) ，从而 \(p_{\mathrm{A}}(x + y) \leqslant p_{\mathrm{A}}(x) + p_{\mathrm{A}}(y)\) 。若 x = y = 0，此不等式仍成立，因为 \(p_{\mathrm{A}}(0) = 0\) 。这就证明了断言 (iii)。

把此不等式应用于 \(x + y\) 与 -y，可得对 E 的点的每个偶 \((x, y)\) ，有

\[
| p _ {\mathrm{A}} (x + y) - p _ {\mathrm{A}} (x) | \leqslant \max (p _ {\mathrm{A}} (y), p _ {\mathrm{A}} (- y)) \leqslant m ^ {- 1} \| y \|.
\]

这就证明映射 \(p_{A}\) 是连续的，引理得证。

映射 \(p_{A}\) 称为凸集 A 的规度。

命题 2. —— 设 E 为 n 维数值空间，A 为 E 的凸紧子集，且 0 是它的一个内点。存在唯一的 E 到自身的双射 u 满足以下三条性质：

\begin{enumerate}
\def\labelenumi{(\roman{enumi})}
\item
  对每个 \(x \in E\) 与每个 \(t \in R_{+}\) ，\(u(tx) = tu(x)\) ；
\item
  对每个 \(x \in E\) ，存在 \(\lambda \in R_{+}\) 使得 \(u(x) = \lambda x\) ；
\item
  我们有 \(u(\mathrm{A}) = \mathbf{B}_n\)
\end{enumerate}

映射 u 是同胚，且通过过渡到子空间，诱导出 A 到 \(B_{n}\) 上的同胚、A 的内部到 \(B_{n}\) 的内部上的同胚，以及 A 在 E 中的边界到球面 \(S_{n-1}\) 上的同胚。

设 \(p_{A}\) 为凸集 A 的规度。设 \(x \in E\) ，\(t \in R_{+}\) ；为使 \(x\) 属于 \(tA\) ，充要条件是 \(p_{\mathrm{A}}(x) \leqslant t\)。由于 A 是紧的，存在实数 M 使得 \(\|x\| \leqslant M\) 对一切 \(x \in A\) 成立。

设 \(u\) 为满足命题条件的映射。我们有 \(u(0) = 0\)。设 \(x \in \mathrm{E}\setminus\{0\}\) ，\(\lambda \in \mathbf{R}_+\) 满足 \(u(x) = \lambda x\)。对一切 \(t \in \mathbf{R}_+\) （使 \(tx \in \mathrm{A}\) ），有 \(u(tx) = t\lambda x\)。由于 \(u\) 是单射，\(\lambda \neq 0\)。由于 \(u(\mathrm{A})\) 包含于 \(\mathbf{B}_n\) ，又有 \(\lambda \leqslant p_{\mathrm{A}}(x)/\|x\|\)。置 \(z = x/\|x\|\) ；它是 \(\mathbf{S}_{n-1}\) 的一点。为使 \(z\) 在 \(u\) 下有 A 中的原像，充要条件是点 \((\lambda \|x\|)^{-1}x\) 属于 A，亦即 \(\lambda \|x\| \geqslant p_{\mathrm{A}}(x)\)，也就是 \(\lambda \geqslant p_{\mathrm{A}}(x)/\|x\|\)。这就蕴含这样的映射 \(u\) 的唯一性。

然后以 \(u\) 表示 E 到自身的、把 0 映到 0 、把 \(x\) 映到 \((p_{\mathrm{A}}(x) / \| x\|)x\) （对每个 \(x \in \mathrm{E}\setminus\{0\}\)）的映射。

由引理 1，\(u\) 在 \(\mathrm{E}\setminus\{0\}\) 的每点处连续。我们有 \(\| u(x)\| = p_{\mathrm{A}}(x)\) （对一切 \(x\in \mathrm{E}\)），且 \(p_{\mathrm{A}}(x)\to 0\) 当 \(x\to 0\) 时（同前引）；因此，\(u\) 在 0 处连续。于是，\(u\) 是连续的。

0 在 u 下唯一的原像是 0。设 \(y \in E\setminus\{0\}\)。为使 E 的元素 x 满足 \(u(x) = y\) ，充要条件是 \(x = (\|y\| / p_{\mathrm{A}}(y))y\)。由此可知 u 是 E 到自身的连续双射。其逆是映射 \(v: y \mapsto (\|y\| / p_{\mathrm{A}}(y))y\)。由于 \(p_{A}\) 连续且仅在 0 处取零值，映射 v 在 \(E\setminus\{0\}\) 的每点处连续；不等式 \(p_{\mathrm{A}}(y) \geqslant \|y\| / \mathrm{M}\) 蕴含 \(\|v(y)\| \leqslant \mathrm{M}\|y\|\) 对一切 \(y \in E\) 成立。由此得出 v 是连续的。于是，映射 u 是 E 到自身的同胚。由于关系 \(p_{\mathrm{A}}(x) \leqslant 1\) 与 \(x \in A\) 等价，我们又有 \(u(\mathrm{A}) = \mathbf{B}_{n}\) 。命题得证。

例. —— 集合 \([0,1]^{n}\) 是 \(R^{n}\) 的内部非空的凸紧子集。由 I, 第 123 页命题 2，它同胚于闭欧氏球。更确切地说，对内部 \(]0,1[^{n}\) 的每个点 x 与开欧氏球的每个点 b，存在 \([0,1]^{n}\) 到 \(B_{n}\) 上的同胚，它把 x 映到 b，并通过过渡到子空间，诱导出 \(]0,1[^{n}\) 到开欧氏球上的同胚，以及 \([0,1]^{n}\) 的边界到球面 \(S_{n-1}\) 上的同胚。

因此，\(R^{n}\) 的每个内部非空的凸紧子集都同胚于一个平行六面体（同前引）。

\subsection{单连通空间}

定义 3. —— 若拓扑空间的一切覆叠都可平凡化，则称它是单连通的。

单连通空间是连通的。事实上，若空间 X 是两个非空开集 U 与 V 的不交并，则 U 到 X 中的典范单射是一个不可平凡化的覆叠。

空空间是单连通的。每个缩为一点的拓扑空间都是单连通的。

注. —— 1) 设 B 为单连通拓扑空间，(E, p) 为 B 的覆叠。若空间 E 连通且非空，则 p 是同胚。例如，设 G 为连通拓扑群，H 为 G 的离散子群，于是典范映射 p: G → G/H 使 G 成为 G/H 的覆叠（I, 第 100 页, 定理 1 的推论 2）。若空间 G/H 单连通，则 p 是同胚且 H 是平凡子群。

\begin{enumerate}
\def\labelenumi{\arabic{enumi})}
\setcounter{enumi}{1}
\tightlist
\item
  设 B 为每个点都有单连通邻域的拓扑空间；则 B 的覆叠的覆叠是 B 的覆叠。事实上，考察 B 的覆叠 (E, p) 与 E 的覆叠 (F, q)。我们来证明 (F, \(p \circ q\)) 是 B 的覆叠。由于问题是 B 上的局部问题，可设空间 B 单连通，从而 E 是 B 的可平凡化覆叠。设 V 为 E 的一个连通分支。它既开又闭，且 \(p \mid V: V \to B\) 是同胚（I, 第 69 页, 命题 1）；因此空间 V 单连通。\(q^{-1}(V)\) 的每个连通分支 W 在 \(q^{-1}(V)\) 中从而在 F 中既开又闭，且映射 q 诱导出 W 到 V 上的同胚。于是映射 \(p \circ q\) 使 F 成为 B 的可平凡化覆叠（同前引）。
\end{enumerate}

命题 3. —— 设 B 为拓扑空间。设 (E, p) 为 B 的覆叠，y 为 E 的一点。设 X 为单连通拓扑空间，f: X → B 为连续映射，x 为 X 的一点且 f(x) = p(y)。存在映射 f 的唯一的连续提升 g: X → E 使得 g(x) = y。

所求的映射 \(g\) 与连续截面 \(s\) 成一一对应（I, 第 9 页, 命题 3），这些截面是覆叠 \(\mathrm{pr}_1: \mathrm{X} \times_{\mathrm{B}} \mathrm{E} \to \mathrm{X}\) 的截面，且满足 \(s(x) = (x, y)\) 。这样的截面存在，因为空间 X 单连通；它是唯一的，因为空间 X 连通（I, 第 34 页, 命题 11 的推论 1）。

例 1. —— 设 X 为单连通拓扑空间，f 为从 X 到 \(C^{*}\) 中的连续函数。回忆（I, 第 101 页, 例 6）映射 \(z \mapsto e^{z}\) 使 C 成为 \(C^{*}\) 的覆叠。由命题 3，存在连续函数 \(h: X \to C\) 使得 \(f(x) = e^{h(x)}\) 对一切 \(x \in X\) 成立。若 \(h': X \to C\) 是另一个连续函数，满足 \(f(x) = e^{h'(x)}\) （对一切 \(x \in X\) ），则存在整数 \(q \in Z\) 使得 \(h' = h + 2\pi iq\)。

类似地，设 \(n\) 为整数 \(>0\) ；映射 \(z \mapsto z^{n}\) 使 \(C^{*}\) 成为自身的覆叠（I, 第 101 页, 例 5）。因此存在连续函数 \(k\) （从 X 到 \(C^{*}\) 中）使得 \(k(x)^{n} = f(x)\) 对一切 \(x \in X\) 成立，例如函数 \(k(x) = e^{h(x)/n}\)。若 \(k' : X \to C^{*}\) 是另一个连续函数，满足 \(k'(x)^{n} = f(x)\) （对一切 \(x \in X\) ），则存在一个 \(n\) 次单位根 \(\mu \in C\) 使得 \(k' = \mu k\)。

推论. —— 设 B 为拓扑空间，b 为 B 的一点。设 E 为 B 的单连通覆叠。对 \(E_{b}\) 的每个点 x，带基点空间 (E, x) 都是 (B, b) 的万有覆叠。

命题 4. —— 两个单连通空间（其中一个是局部连通的）的乘积是单连通空间。

设 X 与 Y 为单连通空间，且设 Y 局部连通。空间 \(X \times Y\) 是连通的，因为 X 与 Y 连通（I, 第 124 页）。设 \((Z, f)\) 为 \(X \times Y\) 的非空覆叠；我们来证明它可平凡化。由 I, 第 70 页推论 2，只需证明对 Z 的每个点 \(z_{0}\) ，存在 f 的连续截面 s 使得 \(s(f(z_{0})) = z_{0}\)。于是设 \(z_{0}\) 为 Z 的一点。置 \((x_{0}, y_{0}) = f(z_{0})\)。X × Y 的子空间 \(X \times \{y_{0}\}\) 同胚于 X。因此，由 Z 限制到 \(X \times \{y_{0}\}\) 上所得的覆叠可平凡化，且有满足 \(\sigma(x_{0}, y_{0}) = z_{0}\) 的连续截面 σ。类似地，对 X 的每个点 x，存在连续截面 \(\tau_{x}\) （f 在 \(\{x\} \times Y\) 上的截面）使得 \(\tau_{x}(x, y_{0}) = \sigma(x, y_{0})\)。对 \((x, y) \in X \times Y\) ，置 \(s(x, y) = \tau_{x}(x, y)\)。映射 s 是 f 的截面，且 \(s(x_{0}, y_{0}) = z_{0}\)。由于空间 Y 连通且局部连通，由 I, 第 35 页定理 1，映射 s 是连续的。

命题 5. —— 连通且局部连通、并且任意两个连通开子集的交都连通的拓扑空间是单连通空间。

设 B 为这样的拓扑空间。设 \((\mathrm{E}, p)\) 为 B 的覆叠，x 为 E 的一点，并记 \(b = p(x)\)。需要证明存在 p 的连续截面 s 使得 \(s(b) = x\) （推论 2, I, 第 70 页）。设 S 为偶 \((\mathrm{U}, s_{\mathrm{U}})\) 组成的集合，其中 U 是 B 的包含 b 的连通开子集，\(s_{U}\) 是 \(p_{\mathrm{U}}: p^{-1}(\mathrm{U}) \to \mathrm{U}\) 的满足 \(s_{\mathrm{U}}(b) = x\) 的连续截面。集合 S 非空（I, 第 34 页, 命题 10）。设 \((\mathrm{U}, s_{\mathrm{U}})\) 与 \((\mathrm{V}, s_{\mathrm{V}})\) 为 S 的元素。此时 \(s_{U} \mid U \cap V\) 与 \(s_{V} \mid U \cap V\) 是 \(p_{U \cap V}\) 的在 b 处取值 x 的连续截面。依假设，\(U \cap V\) 连通；因此 \(s_{U} \mid U \cap V = s_{V} \mid U \cap V\) （I, 第 34 页, 命题 11 的推论 1）。设 A 为当 \((\mathrm{U}, s_{\mathrm{U}})\) 取遍 S 时诸开集 U 的并，并设 s: \(A \to E\) 为满足 \(s \mid U = s_{U}\) （对每个偶 \((\mathrm{U}, s_{\mathrm{U}}) \in S\) ）的唯一映射。集合 A 是开且连通的（TG, I, 第 81 页, 命题 2），它包含 b，且 s 是 \(p_{A}\) 的满足 \(s(b) = x\) 的连续截面。现在只需证明集合 A 是闭的，即可得出它等于 B。

设 \(a\) 为 B 中附着于 A 的点，V 为 \(a\) 的连通开邻域，使得覆叠 E 在 V 上可平凡化。存在一点 \(c\) （属于 \(\mathrm{A} \cap \mathrm{V}\) ）以及连续截面 \(s_{\mathrm{V}}\) （\(p_{\mathrm{V}}\) 的截面），满足 \(s_{\mathrm{V}}(c) = s(c)\) 。设 \(\mathrm{A}'\) 为开集 \(\mathrm{A} \cup \mathrm{V}\)。由于 \(\mathrm{A} \cap \mathrm{V}\) 连通，存在连续截面 \(s'\) （\(p_{\mathrm{A}'}\) 的截面），它同时延拓 \(s\) 与 \(s_{\mathrm{V}}\) （I, 第 35 页, 命题 11 的推论 3）；偶 \((\mathrm{A}', s')\) 属于 \(\mathcal{S}\) ，故 \(\mathrm{A}'\) 包含于 A。于是，\(a\) 属于 A，且 A 是闭的。

推论. —— 实直线 R 的每个区间都是单连通的。

事实上，R 的连通子空间正是区间（TG, IV, 第 8 页, 定理 4），且两个区间的交是区间。

例 2. —— n 维数值空间 \(R^{n}\) （TG, VI, 第 1 页）是单连通的。事实上，它是单连通且局部连通的空间的乘积（I, 第 126 页, 命题 4 及上文命题 5 的推论）。对 \(R^{n}\) 的每个开或闭的平行六面体结论同样成立。平行体与 \(R^{n}\) 中的开或闭欧氏球都是单连通的，因为它们同胚于方块（TG, VI, 第 10 页, 命题 2 及 I, 第 124 页, 例）。

命题 6. —— 设 X 为拓扑空间。设 U₁ 与 U₂ 为 X 的开（相应地，闭）子空间，使得 X = U₁ ∪ U₂。假设 U₁ 与 U₂ 都单连通，且其交 U₁ ∩ U₂ 连通非空。则空间 X 单连通。

设 (E, p) 为 X 的覆叠，y 为 E 的一点。只需证明存在 p 的连续截面 s 使得 \(s(p(y)) = y\) （I, 第 70 页, 命题 1 的推论 2）。例如设 \(p(y)\) 属于 \(U_{1}\)。于是存在连续截面 \(s_{1}: U_{1} \to E\) （\(p_{U_{1}}\) 的截面）使得 \(s_{1}(p(y)) = y\)。设 x 为 \(U_{1} \cap U_{2}\) 的一点；存在连续截面 \(s_{2}\) （\(p_{U_{2}}\) 的截面）使得 \(s_{2}(x) = s_{1}(x)\)。由 I, 第 34 页命题 11 的推论 3，存在同时延拓 \(s_{1}\) 与 \(s_{2}\) 的 p 的连续截面 s。

例. —— 3) 对 \(n \geqslant 2\)，球面 \(\mathbf{S}_{n}\) （TG, VI, 第 10 页）是单连通的。事实上，球面 \(\mathbf{S}_{n}\) 是两个同胚于 \(\mathbf{B}_{n}\) 的闭半球之并，这两个闭半球的交同胚于

\(S_{n-1}\) （参见 TG, VI, 第 12 页）。对 \(n \geqslant 2\)，球面 \(S_{n-1}\) 连通，断言得证。

另一方面，圆 \(S_{1}\) 不是单连通的。事实上，连续映射 \(p: R \to S_{1}\) 由 \(p(\theta) = (\cos \theta, \sin \theta)\) 所定义，它使 R 成为 \(S_{1}\) 的无限次覆叠，它是连通的，因而不可平凡化。

\begin{enumerate}
\def\labelenumi{\arabic{enumi})}
\setcounter{enumi}{3}
\tightlist
\item
  设 E 为 R 上有限维 n 维向量空间，F 为 E 的余维数为 \(p \geqslant 3\) 的仿射子空间。集合 \(R^{p}\setminus\{0\}\) 同胚于 \(R \times S_{p-1}\) （TG, VI, 第 10 页, 推论 2），故它是单连通的，因为 R 与 \(S_{p-1}\) 都局部连通且单连通（I, 第 126 页, 命题 4）。集合 \(E\setminus F\) 同胚于 \(R^{n-p} \times (R^{p}\setminus\{0\})\) ，因而是单连通的（同前引）。
\end{enumerate}

命题 7. —— 设 X 为拓扑空间。设 \(U_{1}\) 与 \(U_{2}\) 为 X 的连通开（相应地，闭）子空间，使得 \(X = U_{1} \cup U_{2}\)。若空间 X 单连通，则 \(U_{1} \cap U_{2}\) 连通。

置 \(U = U_{1} \cap U_{2}\) ，并反设空间 U 不连通。我们将构造 X 的一个无限次数的连通覆叠；这样的覆叠不可平凡化。

依假设，集合 U 是两个不交非空、且在 X 中为开（相应地，闭）的集合 A 与 B 的并。对 \(i \in \{1, 2\}\) ，设 \(Y_i\) 为 \(U_i\)-空间 \((\mathrm{U}_i \times \mathbf{Z}, \mathrm{pr}_1)\) ，\(Z_i\) 为子空间 \(U \times \mathbf{Z}\) （\(Y_i\) 的子空间）。映射 \(h: Z_1 \to Z_2\) 由

\[
h (x, n) = \left\{ \begin{array}{l l} (x, n) & \text {if } x\in A \text { and } n\in \mathbf {Z} \\ (x, n + 1) & \text {if } x\in B \text { and } n\in \mathbf {Z} \end{array} \right.
\]

所定义，它是同胚。设 Y 为把 \(Y_{1}\) 与 \(Y_{2}\) 沿 \(Z_{1}\) 与 \(Z_{2}\) 借助同胚 h 粘合起来所得的空间（TG, I, 第 17 页）。

对 \(i \in \{1,2\}\) ，典范映射 \(g_i\) （从 \(Y_i\) 到 Y 中）是 \(Y_i\) 到 Y 的一个开（相应地，闭）子集上的同胚（同前引, 命题 9）。对每个整数 \(n \in \mathbf{Z}\) ，集合 \(g_i(U_i \times \{n\})\) （\(i \in \{1,2\}\)）都是连通的；由于 A 与 B 非空，\(g_1(U_1 \times \{n\})\) 与 \(g_2(U_2 \times \{n\})\) 以及 \(g_2(U_2 \times \{n+1\})\) 相交。由此得出空间 Y 连通（TG, I, 第 81 页, 推论）。

对 \(x \in \mathrm{U}\) 与 \(n \in \mathbf{Z}\) ，有 \((\mathrm{pr}_1 \circ h)(x, n) = x = \mathrm{pr}_1(x, n)\)。因此存在唯一的连续映射 \(p \colon \mathrm{Y} \to \mathrm{X}\) 使得 \(p \circ g_i = \mathrm{pr}_1\) （对 \(i \in \{1, 2\}\)）。我们来证明 X-空间 \((\mathrm{Y}, p)\) 是一个覆叠。

按构造，映射 p 的纤维同胚于离散空间 Z。

对 \(i \in \{1,2\}\) ，映射 \(g_i\) 通过过渡到子空间，定义出 \(U_i\)-空间之间从 \(\mathrm{U}_i \times \mathbf{Z}\) 到 \(p^{-1}(\mathrm{U}_i)\) 上的一个同构。

按空间 Y 的定义，存在唯一的从 \((\mathrm{X}\setminus\mathrm{A})\times\mathbf{Z}\) 到 Y 中的映射 k，使得 \(k(x,n)=g_{1}(x,n)\) （当 \(x\in(\mathrm{U}_{1}\setminus\mathrm{A})\times\mathbf{Z}\) ），\(k(x,n)=g_{2}(x,n-1)\) （当 \(x\in(\mathrm{U}_{2}\setminus\mathrm{A})\times\mathbf{Z}\) ）；它是 \((\mathrm{X}\setminus\mathrm{A})\)-空间之间从 \((\mathrm{X}\setminus\mathrm{A})\times\mathbf{Z}\) 到 \(p^{-1}(\mathrm{X}\setminus\mathrm{A})\) 上的同构。同样，存在唯一的映射 \(k'\) （从 \((\mathrm{X}\setminus\mathrm{B})\times\mathbf{Z}\) 到 Y 中），它与 \(g_{1}\) 在 \((\mathrm{U}_{1}\setminus\mathrm{B})\times\mathbf{Z}\) 上相合，与 \(g_{2}\) 在 \((\mathrm{U}_{2}\setminus\mathrm{B})\times\mathbf{Z}\) 上相合；它是 \((\mathrm{X}\setminus\mathrm{B})\)-空间之间从 \((\mathrm{X}\setminus\mathrm{B})\times\mathbf{Z}\) 到 \(p^{-1}(\mathrm{X}\setminus\mathrm{B})\) 上的同构。

这就证明 X-空间 \((Y,p)\) 在部分 \(U_{1}\)、\(U_{2}\)、\(X\setminus A\) 与 \(X\setminus B\) 上都可平凡化。我们有 \(U_{1} \cup U_{2} = X\) ；若 \(U_{1}\) 与 \(U_{2}\) 在 X 中是开的，这就证明 \((Y,p)\) 是 X 的覆叠。当 \(U_{1}\) 与 \(U_{2}\) 在 X 中是闭时结论同样成立，因为此时 \(X\setminus A\) 与 \(X\setminus B\) 是 X 的开子集且其并为 X。命题至此得证。

推论. —— 设 X 为单连通拓扑空间，A 为 X 的连通子集。若 A 的余集连通，则其边界也连通。

设 \(X_{1}\) 与 \(X_{2}\) 分别为 A 与 \(X\setminus A\) 的闭包。集合 \(X_{1}\) 与 \(X_{2}\) 是闭且连通的（TG, I, 第 81 页, 命题 1）；我们有 \(X_{1} \cup X_{2} = X\) ，且其交 \(X_{1} \cap X_{2}\) 等于 A 的边界。于是只需应用命题 7。
